%! Author = crist
%! Date = 23/10/2026
%
% Nomenclatura de Química Inorgánica — revisión según:
%   · RSEQ, "Recomendaciones para la enseñanza de la nomenclatura de Química
%     Inorgánica en la Enseñanza Secundaria" (programación en espiral).
%   · IUPAC 2005 (Libro Rojo) y la lista de nombres vulgares admitidos.
% Las figuras y esquemas se generan desde images/html/figuras_inorganica.html
%   (?solo=<id> muestra una sola figura) y se exportan en vectorial a
%   images/pdf/<id>.pdf (los que usa LaTeX) e images/svg/<id>.svg.
% La versión anterior de este archivo está en inorganica_nuevo_ORIGINAL_backup.tex

% Preamble
\documentclass[12pt]{article}
% Packages
\usepackage[english,spanish, es-tabla]{babel}
\usepackage[top=0.5cm, left=1.5cm, right=1.5cm, bottom=1.5cm]{geometry}
\usepackage{estilocrispropio}
\usepackage{longtable} % tablas largas con formato booktabs
% Tipos de columna de párrafo (los mismos que LegrandOrangeBookES) para las tablas booktabs
\newcolumntype{R}[1]{>{\raggedleft\arraybackslash}p{#1}}
\newcolumntype{L}[1]{>{\raggedright\arraybackslash}p{#1}}
\newcolumntype{C}[1]{>{\centering\arraybackslash}p{#1}}


% CARGA DE NEWPX CON SUS OPCIONES PRINCIPALES
\usepackage[
    scale=.98,        % Escala de la fuente (por defecto 1) [2, 5]
    osf,              % Cifras de estilo antiguo (Old-style figures) en texto [2, 6]
    p,                % Cifras proporcionales en texto [2, 6]
    theoremfont,      % Fuente especial para teoremas (itálicas con puntuación recta) [3, 7]
    trueslanted,      % \textsl produce texto inclinado real en lugar de itálicas [3, 7]
    largesc,          % Aumenta el tamaño de las versalitas un 5.33% [3, 8]
    tighter,          % Ajusta el espaciado entre palabras para que sea más compacto [3, 9]
    vvarbb,           % Usa el estilo STIX para la negrita de pizarra (\mathbb) [10-12]
    varg,             % Usa una 'g' itálica matemática de un solo lazo [13, 14]
    boldermath,       % Hace que las itálicas romanas en matemáticas sean más gruesas [15]
    amsthm            % Carga soporte para amsthm automáticamente [10]
]{newpx}

%%%%%%%%%%%%%%%%%%%%%%%%%%%%
% Ejercicios y soluciones (paquete answers, cargado en estilocrispropio)
%   \begin{exercise} ... \begin{solu} ... \end{solu} \end{exercise}
%   Las soluciones se escriben en ans1.tex y se imprimen al final.
%   (Se quita xsim porque también define el entorno exercise.)
%%%%%%%%%%%%%%%%%%%%%%%%%%%
\theoremstyle{mygreennumbox}
\newtheorem{ejercicioT}{Ejercicio}[section]
\newenvironment{exercise}{\begin{eBoxmygreen}\begin{ejercicioT}}{\hfill{\color{mygreen}\tiny\ensuremath{\blacksquare}}\end{ejercicioT}\end{eBoxmygreen}}
\renewenvironment{Solution}[1]{%
    \begin{eBoxmyBlue}%
    \noindent{\small\sffamily\bfseries\color{myBlue}Ejercicio~#1}\par\smallskip%
}{%
    \end{eBoxmyBlue}%
}

%%%%%%%%%%%%%%%%%%%%%%%%%%%%
% Numeraciones dentro de las cajas de color: las bolas numeradas toman el
% color del borde de la caja (vale para \numberballgreen, \numberballblue,
% \numberballred e \itembolascolor). Fuera de las cajas no cambia nada.
%%%%%%%%%%%%%%%%%%%%%%%%%%%%
\newcommand{\numeracioncolor}[1]{%
    \colorlet{colornumeracion}{#1}%
    \renewcommand*{\numberballgreen}{\protect\cnumero{\arabic*}{colornumeracion}}%
    \renewcommand*{\numberballblue}{\protect\cnumero{\arabic*}{colornumeracion}}%
    \renewcommand*{\numberballred}{\protect\cnumero{\arabic*}{colornumeracion}}%
    \renewcommand*{\itembolascolor}[2]{\protect\cnumero{##1}{colornumeracion}}%
}
\AtBeginEnvironment{consejos}{\numeracioncolor{yellow!75!black}}
\AtBeginEnvironment{ejemplo}{\numeracioncolor{green!65!black}}
\AtBeginEnvironment{ejercicios}{\numeracioncolor{green!65!black}}
\AtBeginEnvironment{definicion}{\numeracioncolor{red!75!black}}
\AtBeginEnvironment{solucion}{\numeracioncolor{blue!90}}
\AtBeginEnvironment{marker}{\numeracioncolor{yellow!50!black}}
\AtBeginEnvironment{exercise}{\numeracioncolor{mygreen}}
\AtBeginEnvironment{Solution}{\numeracioncolor{myBlue}}

\sisetup{%
%alsoload = synchem,
    detect-family,%
    per-mode = fraction,
    display-per-mode = fraction,
%inline-per-mode = symbol,
%per-mode = symbol-or-fraction,
    output-decimal-marker = {,},
    inter-unit-product = \cdot,
    exponent-product = \cdot,
    retain-explicit-plus % consigo que si pongo + en el número aparezca
}

% --- Comandos propios de este documento ---
% Esquema de una regla de nomenclatura (figura vectorial generada desde HTML)
\newcommand{\esquemafig}[2][0.85]{\begin{center}\includegraphics[width=#1\textwidth]{images/pdf/#2.pdf}\end{center}}
% Cabecera de los apartados "¿Cómo se formulan / nombran?"
\newcommand{\formulan}{\vspace*{2mm}\noindent\xmybox{\bfseries ¿Cómo se formulan?}\par\smallskip}
\newcommand{\nombran}{\vspace*{2mm}\noindent\xmybox{\bfseries ¿Cómo se nombran?}\par\smallskip}
\newcommand{\ejemplos}{\par\noindent\textbf{\uline{Ejemplos}}:\par}
% Nombre incorrecto (tachado en rojo)
\newcommand{\mal}[1]{\textcolor{red}{\sout{#1}}}
% Marca de contenido de ampliación (icono + etiqueta); \ampliacionicono solo el icono
\newcommand{\ampliacionicono}{\texorpdfstring{\textcolor{myOrange3}{\faRocket}}{}}
\newcommand{\ampliacion}{\texorpdfstring{\tcbox[on line, colback=myOrange3!12, colframe=myOrange3, boxrule=0.6pt, arc=2pt, boxsep=0pt, left=3pt, right=3pt, top=1pt, bottom=1pt]{\footnotesize\sffamily\bfseries\textcolor{myOrange3}{\faRocket\ Ampliación}}}{(Ampliación)}}
% Etiquetas que indican con qué se nombra (teoría, enunciados y soluciones)
\definecolor{colornumox}{HTML}{7A2BB5}
\newcommand{\etiquetanom}[2]{\tcbox[on line, colback=#1!10, colframe=#1, boxrule=0.5pt, arc=2pt, boxsep=0pt, left=2pt, right=2pt, top=0.5pt, bottom=0.5pt]{\scriptsize\sffamily\bfseries\textcolor{#1}{#2}}}
\newcommand{\conpref}{\etiquetanom{myBlue}{prefijos}}
\newcommand{\conox}{\etiquetanom{colornumox}{nº de oxidación}}
\newcommand{\concarga}{\etiquetanom{myOrange3}{nº de carga}}
% Figura generada desde HTML (PDF vectorial; la versión SVG está en images/svg)
\newcommand{\figura}[2][0.95]{\begin{figure}[H]\centering\includegraphics[width=#1\textwidth]{images/pdf/#2.pdf}\end{figure}}

% Document
\begin{document}
\begin{titlepage}
    \vspace*{5cm}
    \begin{center}
        {\Huge{ \bfseries Nomenclatura Inorgánica}}
    \end{center}
    \vspace*{2cm}
    \begin{center}
        {\huge  IUPAC 2005}\par\medskip
        {\Large Adaptada a las recomendaciones de la RSEQ para Secundaria y Bachillerato}
    \end{center}
    \vspace*{5cm}
    \begin{center}
        {\Large IES José Conde García (Almansa)}
    \end{center}
\end{titlepage}

\tableofcontents

\setstretch{1.3}
\newpage

% Abrimos el fichero en el que se guardan las soluciones de los ejercicios
\Opensolutionfile{ans}[ans1]

%=====================================================================
\section{Introducción}
%=====================================================================

La \negrita{nomenclatura química} es el conjunto de reglas que sirven para asignar \negrita{descriptores} (nombres y fórmulas) a las sustancias químicas, de forma que cualquier persona del mundo sepa de qué sustancia estamos hablando. Nombrar (de la fórmula al nombre) y formular (del nombre a la fórmula) son, por tanto, dos caras de la misma moneda: la nomenclatura.

A continuación, se dan una serie de conceptos que es interesante conocer:
\begin{enumerate}[label=\numberballred]
    \item \negrita{Tabla periódica}: disposición de los elementos químicos ordenados por su número atómico creciente en columnas (\negrita{grupos}) y filas (\negrita{periodos}), de modo que los elementos de un mismo grupo tienen propiedades parecidas.
    \item \negrita{Fórmula}: combinación de símbolos químicos y subíndices que expresa la composición de una sustancia, es decir, \negrita{qué} elementos la forman y en \negrita{qué proporción}. Por ejemplo, \ch{H2O} indica que hay dos átomos de hidrógeno por cada átomo de oxígeno.
    \item \negrita{Número (o estado) de oxidación}: carga que tendría un átomo si todos sus enlaces fueran iónicos. Puede ser positivo, negativo o cero y se escribe con el signo \negrita{delante}: \(+1\), \(+2\), \(0\), \(-1\), \(-2\)\dots
    \item \negrita{Número de carga}: carga real de un ion. Se escribe con el signo \negrita{detrás}: \ch{Fe^3+} tiene número de carga \(3+\).
    \item \negrita{Valencia}: número de enlaces que forma un átomo. Es un concepto distinto del número de oxidación y \negrita{no} lo usaremos para nombrar.
\end{enumerate}

Antes de empezar, conviene tener una visión general de los tipos de compuestos que vamos a estudiar. La primera pregunta que debes hacerte siempre ante una fórmula es: \emph{¿cuántos elementos tiene y cuáles son?}

\begin{marker}
    Los apartados marcados con \ampliacion{} (o con el icono \ampliacionicono{} en el índice) son de \negrita{ampliación}: amplían lo básico y no es necesario estudiarlos la primera vez que se ve el tema.
\end{marker}

\figura{fig-mapa}

%=====================================================================
\section{Herramientas básicas}
%=====================================================================

\subsection{¿Qué elemento se escribe a la derecha? La secuencia de los elementos}\label{sec:secuencia}

En un compuesto binario siempre hay un constituyente que se escribe a la izquierda y otro a la derecha. La IUPAC fija el orden mediante una \negrita{secuencia de los elementos} (Tabla VI del Libro Rojo). Esta tabla \negrita{no} es una ordenación por electronegatividades: por ejemplo, el oxígeno es más electronegativo que el cloro y, sin embargo, en la secuencia el cloro aparece antes que el oxígeno. Por eso es mejor hablar simplemente de \negrita{orden de colocación}.

\figura{fig-secuencia}

\begin{marker}
    Consecuencia importante (IUPAC 2005): los halógenos (F, Cl, Br, I) están en la secuencia a la derecha del oxígeno. Por eso sus compuestos con oxígeno \negrita{no son óxidos}, sino halogenuros de oxígeno: \ch{OCl2} es dicloruro de oxígeno (y no \mal{\ch{Cl2O}}, \mal{óxido de dicloro}) y \ch{OF2} es difluoruro de oxígeno.
\end{marker}

\subsection{Números de oxidación que hay que conocer}\label{sec:eo}

No es necesario memorizar largas tablas de números de oxidación. Basta con conocer unos pocos, apoyándonos en la tabla periódica; el resto se deduce con la regla de la suma cero (apartado~\ref{sec:sumacero}).

\figura{fig-tabla-eo}

\subsubsection{Reglas generales según el grupo}\label{sec:eo_grupos}

En los elementos de los \negrita{grupos principales} (1, 2 y 13 al 17) no hace falta memorizar los números de oxidación: se deducen de su \negrita{posición en la tabla periódica}.

\figura{fig-eo-grupos}

\begin{enumerate}[label=\numberballblue]
    \item \negrita{Elemento sin combinar}: su número de oxidación es \(0\) (\ch{Na}, \ch{Fe}, \ch{O2}, \ch{Cl2}\dots).
    \item \negrita{Metales de los grupos 1, 2 y 13}: solo tienen un número, positivo e igual a la \negrita{última cifra del grupo}. Grupo 1 \(\to +1\); grupo 2 \(\to +2\); Al (grupo 13) \(\to +3\).
    \item \negrita{Mayor número positivo}: es la \negrita{última cifra del grupo}. C (14) \(\to +4\); N (15) \(\to +5\); S (16) \(\to +6\); Cl (17) \(\to +7\). Los no metales lo usan, sobre todo, cuando van con el oxígeno (\ch{CO2}, \ch{SO3}).
    \item \negrita{Número negativo de los no metales}: es \negrita{grupo \(-\) 18}. C \(\to -4\); N \(\to -3\); S \(\to -2\); Cl \(\to -1\). Lo usan cuando se escriben a la \negrita{derecha} de la fórmula, con metales o con hidrógeno: \ch{NaCl} (Cl \(-1\)), \ch{H2S} (S \(-2\)), \ch{NH3} (N \(-3\)), \ch{CH4} (C \(-4\)).
    \item \negrita{Los demás números positivos}: se obtienen \negrita{restando de 2 en 2} desde el mayor. Así, los grupos \negrita{impares} tienen números \negrita{impares} y los \negrita{pares}, números \negrita{pares}:
          \begin{center}
              Cl: \(+7,\ +5,\ +3,\ +1\) \qquad S: \(+6,\ +4,\ +2\) \qquad N: \(+5,\ +3,\ +1\) \qquad C: \(+4,\ +2\)
          \end{center}
\end{enumerate}

\begin{marker}
    \negrita{Excepciones que hay que recordar}:
    \begin{itemize}
        \item \negrita{Hidrógeno}: \(+1\) con los no metales y \(-1\) con los metales (hidruros).
        \item \negrita{Oxígeno}: casi siempre \(-2\); en los peróxidos \(-1\) y, con los halógenos, positivo (\ch{OF2}: O \(+2\)).
        \item \negrita{Flúor}: solo \(-1\), porque es el elemento más electronegativo.
        \item El \negrita{nitrógeno} también puede actuar con \(+2\) y \(+4\) (\ch{NO}, \ch{NO2}).
    \end{itemize}
\end{marker}

\negrita{Metales de transición}. Para ellos no hay una regla sencilla. Los más habituales son:

\begin{center}
    \SetTblrInner{rowsep=3pt}
    % --- tabla antigua ---
    % \begin{tblr}{colspec={Q[l,m] Q[c,m] Q[l,m] Q[c,m]}, row{1}={bg=myOrange3!90!white, fg=white, font=\sffamily}, hline{1,2,Z}={1pt}, hline{3-Y}={dashed}}
        % Elemento & Nº de oxidación & Elemento & Nº de oxidación \\
        % Ag & \(+1\) & Zn, Cd & \(+2\) \\
        % Cu, Hg & \(+1,\ +2\) & Au & \(+1,\ +3\) \\
        % Fe, Co, Ni & \(+2,\ +3\) & Sn, Pb (grupo 14) & \(+2,\ +4\) \\
        % Cr & \(+2,\ +3,\ +6\) & Mn & \(+2,\ +3,\ +4,\ +6,\ +7\) \\
    % \end{tblr}
    % --- tabla nueva ---
    \begin{tabular}{@{}lclc@{}}
        \toprule
        \textbf{Elemento} & \textbf{Nº de oxidación} & \textbf{Elemento} & \textbf{Nº de oxidación} \\
        \midrule
        Ag & \(+1\) & Zn, Cd & \(+2\) \\[3pt]
        Cu, Hg & \(+1,\ +2\) & Au & \(+1,\ +3\) \\[3pt]
        Fe, Co, Ni & \(+2,\ +3\) & Sn, Pb (grupo 14) & \(+2,\ +4\) \\[3pt]
        Cr & \(+2,\ +3,\ +6\) & Mn & \(+2,\ +3,\ +4,\ +6,\ +7\) \\
        \bottomrule
    \end{tabular}
\end{center}

En los ejercicios no hace falta saberlos de memoria: el número de oxidación del metal se deduce de la fórmula con la regla de la suma cero (apartado~\ref{sec:sumacero}).

\begin{exercise}
    \textbf{1.} Usando las reglas del grupo, indica el mayor número de oxidación positivo y el negativo de: \textbf{a)} P; \textbf{b)} Se; \textbf{c)} Br; \textbf{d)} Si; \textbf{e)} Ba; \textbf{f)} Al.\\
    \textbf{2.} ¿Qué números de oxidación positivos pueden tener el yodo y el azufre?
    \begin{solu}
        \textcolor{blue}{Mayor positivo = última cifra del grupo; negativo = grupo \(-\) 18 (solo los no metales).}\\
        \textbf{1. a)} P (grupo 15): \textcolor{red}{\(+5\) y \(-3\)}; \textbf{b)} Se (16): \textcolor{red}{\(+6\) y \(-2\)}; \textbf{c)} Br (17): \textcolor{red}{\(+7\) y \(-1\)}; \textbf{d)} Si (14): \textcolor{red}{\(+4\) y \(-4\)}; \textbf{e)} Ba (2): \textcolor{red}{\(+2\)}, no tiene negativo (es un metal); \textbf{f)} Al (13): \textcolor{red}{\(+3\)}, no tiene negativo\\
        \textbf{2.} \textcolor{blue}{Restamos de 2 en 2 desde el mayor.} Yodo (17): \textcolor{red}{\(+7,\ +5,\ +3,\ +1\)}; azufre (16): \textcolor{red}{\(+6,\ +4,\ +2\)}
    \end{solu}
\end{exercise}

\subsection{Número de oxidación y número de carga: no son lo mismo}\label{sec:noxcarga}

Son dos ideas distintas que a menudo se confunden. Las dos sirven para indicar la proporción de los elementos en un nombre, pero no se usan en los mismos compuestos ni se escriben igual.

\begin{definicion}{número de oxidación \conox}
    \begin{itemize}
        \item \negrita{Qué es}: la carga que \emph{tendría} un átomo si todos sus enlaces fueran iónicos. Es un número que se asigna a cada \negrita{átomo}, no a los grupos de átomos.
        \item \negrita{Dónde se usa}: en \negrita{todos} los compuestos, tanto iónicos (\ch{Fe2O3}) como covalentes (\ch{SO3}, \ch{PCl5}).
        \item \negrita{En los cálculos y sobre la fórmula}: el signo va \negrita{delante} del número: \ox{+3,Fe}, \ox{-2,O}.
        \item \negrita{En el nombre}: en \negrita{números romanos}, entre paréntesis, sin signo y sin espacio: óxido de hierro(III), óxido de azufre(VI).
    \end{itemize}
\end{definicion}

\begin{definicion}{número de carga \concarga}
    \begin{itemize}
        \item \negrita{Qué es}: la carga \emph{real} de un \negrita{ion}, ya sea un átomo (\ch{Fe^3+}) o un grupo de átomos (\ch{SO4^2-}).
        \item \negrita{Dónde se usa}: \negrita{solo} en los iones y en los compuestos iónicos (metal + no metal, hidróxidos, oxosales). \negrita{Nunca} en compuestos covalentes como \ch{SO3}, \ch{CO2} o \ch{PCl5}, porque no tienen iones.
        \item \negrita{En la fórmula}: arriba a la derecha, el signo va \negrita{detrás} del número: \ch{Fe^3+}, \ch{SO4^2-}. El 1 no se escribe: \ch{Na+}, \ch{Cl-}.
        \item \negrita{En el nombre}: en \negrita{números arábigos} seguidos del signo, entre paréntesis y sin espacio: óxido de hierro(3+), hierro(3+).
        \item Los grupos de átomos con carga, como el hidróxido \ch{OH-}, tienen \negrita{número de carga} (\(1-\)), no número de oxidación.
    \end{itemize}
\end{definicion}

\begin{center}
    \SetTblrInner{rowsep=4pt}
    % --- tabla antigua ---
    % \begin{longtblr}[caption={Número de oxidación frente a número de carga}]
        % {colspec={Q[l,m,3.2cm] X[l,m] X[l,m]}, row{1}={bg=myOrange3!90!white, fg=white, font=\sffamily}, rowhead=1}
        % \hline
        % & Número de oxidación \conox & Número de carga \concarga \\
        % \hline
        % ¿Qué indica? & Carga que tendría el átomo si los enlaces fueran iónicos & Carga real del ion \\
        % ¿Dónde se usa? & En todos los compuestos & Solo en iones y compuestos iónicos \\
        % En la fórmula & Signo delante: \(+3\) & Signo detrás: \ch{Fe^3+} \\
        % En el nombre & Romanos: hierro(III) & Arábigos y signo: hierro(3+) \\
        % \ch{Fe2O3} (iónico) & óxido de hierro(III) & óxido de hierro(3+) \\
        % \ch{SO3} (covalente) & óxido de azufre(VI) & \mal{óxido de azufre(6+)}: no se puede \\
        % \hline
    % \end{longtblr}
    % --- tabla nueva ---
    \begin{longtable}{@{}L{3.2cm} L{0.38\linewidth} L{0.38\linewidth}@{}}
        \caption{Número de oxidación frente a número de carga} \\
        \toprule
         & \textbf{Número de oxidación \conox} & \textbf{Número de carga \concarga} \\
        \midrule
        \endfirsthead
        \toprule
         & \textbf{Número de oxidación \conox} & \textbf{Número de carga \concarga} \\
        \midrule
        \endhead
        \bottomrule
        \endlastfoot
        ¿Qué indica? & Carga que tendría el átomo si los enlaces fueran iónicos & Carga real del ion \\[3pt]
        ¿Dónde se usa? & En todos los compuestos & Solo en iones y compuestos iónicos \\[3pt]
        En la fórmula & Signo delante: \(+3\) & Signo detrás: \ch{Fe^3+} \\[3pt]
        En el nombre & Romanos: hierro(III) & Arábigos y signo: hierro(3+) \\[3pt]
        \ch{Fe2O3} (iónico) & óxido de hierro(III) & óxido de hierro(3+) \\[3pt]
        \ch{SO3} (covalente) & óxido de azufre(VI) & \mal{óxido de azufre(6+)}: no se puede \\
    \end{longtable}
\end{center}

\begin{marker}
    \negrita{Fíjate en lo que pide el enunciado.} En los ejercicios se indica con una etiqueta cómo hay que nombrar: \conpref{} con prefijos multiplicadores, \conox{} con el número de oxidación (romanos) o \concarga{} con el número de carga (arábigos y signo, solo compuestos iónicos). Las soluciones usan las mismas etiquetas.
\end{marker}

\begin{exercise}
    Para cada sustancia, indica el \conox{} del elemento destacado y, si es un compuesto iónico, el \concarga{} de su ion: \textbf{a)} Fe en \ch{FeCl3}; \textbf{b)} S en \ch{SO2}; \textbf{c)} Cu en \ch{CuO}; \textbf{d)} P en \ch{PCl3}; \textbf{e)} Ca en \ch{CaF2}; \textbf{f)} el grupo \ch{OH} en \ch{NaOH}.
    \begin{solu}
        \textcolor{blue}{El número de carga solo existe si el compuesto es iónico (metal + no metal).}\\
        \textbf{a)} \conox{} \textcolor{red}{\(+3\)} \quad \concarga{} \textcolor{red}{\(3+\)} (ion \ch{Fe^3+})\\
        \textbf{b)} \conox{} \textcolor{red}{\(+4\)} \quad \concarga{} \textcolor{red}{no tiene}: el \ch{SO2} es covalente\\
        \textbf{c)} \conox{} \textcolor{red}{\(+2\)} \quad \concarga{} \textcolor{red}{\(2+\)} (ion \ch{Cu^2+})\\
        \textbf{d)} \conox{} \textcolor{red}{\(+3\)} \quad \concarga{} \textcolor{red}{no tiene}: el \ch{PCl3} es covalente\\
        \textbf{e)} \conox{} \textcolor{red}{\(+2\)} \quad \concarga{} \textcolor{red}{\(2+\)} (ion \ch{Ca^2+})\\
        \textbf{f)} \conox{} \textcolor{red}{no se asigna a un grupo} (sí a cada átomo: O \(-2\), H \(+1\)) \quad \concarga{} \textcolor{red}{\(1-\)} (ion \ch{OH-})
    \end{solu}
\end{exercise}

\begin{exercise}
    Corrige los errores: \textbf{a)} cloruro de hierro(+3); \textbf{b)} \ch{Ca^{+2}}; \textbf{c)} óxido de azufre(6+); \textbf{d)} cloruro de hierro (III); \textbf{e)} hierro(3).
    \begin{solu}
        \textbf{a)} \textcolor{red}{cloruro de hierro(3+)} \concarga{} o \textcolor{red}{cloruro de hierro(III)} \conox\\
        \textbf{b)} \textcolor{red}{\ch{Ca^2+}}: en la carga el signo va detrás del número\\
        \textbf{c)} \textcolor{red}{óxido de azufre(VI)} \conox: el \ch{SO3} es covalente y no admite número de carga\\
        \textbf{d)} \textcolor{red}{cloruro de hierro(III)}: sin espacio antes del paréntesis\\
        \textbf{e)} \textcolor{red}{hierro(3+)} \concarga: el número de carga lleva siempre el signo
    \end{solu}
\end{exercise}

\subsection{La regla de la suma cero}\label{sec:sumacero}

\begin{impt}
    \centering En un compuesto neutro, la suma de los números de oxidación de todos sus átomos es \textbf{cero}.\\ En un ion, la suma es igual a la carga del ion.
\end{impt}

Esta regla sirve para dos cosas: escribir la fórmula a partir del nombre y averiguar el número de oxidación de un elemento a partir de la fórmula. Es más segura que el método de «cruzar los números y simplificar», que conduce a errores en muchos casos (por ejemplo, el tetraóxido de dinitrógeno, \ch{N2O4}, no se simplifica).

\figura{fig-suma-cero}

\begin{exercise}
    Calcula, con la regla de la suma cero, el \conox{} del elemento destacado: \textbf{a)} Cr en \ch{Cr2O3}; \textbf{b)} S en \ch{SO3}; \textbf{c)} N en \ch{N2O5}; \textbf{d)} Mn en \ch{MnO2}; \textbf{e)} Cl en \ch{HClO4}.
    \begin{solu}
        \textcolor{blue}{El O es \(-2\) y el H es \(+1\); la suma debe ser cero.}\\
        \textbf{a)} \(2x + 3\cdot(-2) = 0 \Rightarrow\) \textcolor{red}{\(x = +3\)}\\
        \textbf{b)} \(x + 3\cdot(-2) = 0 \Rightarrow\) \textcolor{red}{\(x = +6\)}\\
        \textbf{c)} \(2x + 5\cdot(-2) = 0 \Rightarrow\) \textcolor{red}{\(x = +5\)}\\
        \textbf{d)} \(x + 2\cdot(-2) = 0 \Rightarrow\) \textcolor{red}{\(x = +4\)}\\
        \textbf{e)} \((+1) + x + 4\cdot(-2) = 0 \Rightarrow\) \textcolor{red}{\(x = +7\)}
    \end{solu}
\end{exercise}

\subsection{Prefijos multiplicadores}\label{sec:prefijos}

\figura[0.85]{fig-prefijos}

\begin{marker}
    El prefijo \blueremark{mono} solo se utiliza cuando hace falta distinguir entre varios compuestos de los mismos elementos: \ch{CO} \negrita{monóxido} de carbono (porque existe el \ch{CO2}), \ch{FeO} monóxido de hierro. En cambio, \ch{Na2O} es \negrita{óxido de disodio} (no \mal{monóxido de disodio}).
\end{marker}

\subsection{Tipos de nomenclatura}

Según la IUPAC, los sistemas de nomenclatura de química inorgánica son:
\begin{enumerate}[label=\numberballblue]
    \item \negrita{Nomenclatura de composición} (o estequiométrica). Solo informa de qué elementos hay y en qué proporción. Es la que usaremos para los compuestos binarios y los hidróxidos. La proporción se puede indicar de tres formas:
          \begin{enumerate}[label=\itembolascolor{\arabic*}{blue}]
              \item Con \negrita{prefijos multiplicadores}: \ch{Fe2O3} trióxido de dihierro. \textcolor{red}{Es la forma preferida y la primera que se aprende.}
              \item Con el \negrita{número de oxidación} en números romanos: \ch{Fe2O3} óxido de hierro(III).
              \item Con el \negrita{número de carga} en números arábigos (solo compuestos iónicos): \ch{Fe2O3} óxido de hierro(3+).
          \end{enumerate}
    \item \negrita{Nomenclatura de sustitución}. Se basa en unos \negrita{hidruros progenitores} (\ch{CH4} metano, \ch{NH3} azano, \ch{PH3} fosfano\dots) a partir de los cuales se nombran otros compuestos sustituyendo hidrógenos. Es la base de la nomenclatura orgánica.
    \item \negrita{Nomenclatura de adición}. Considera que el compuesto está formado por un átomo central al que se unen otros átomos o grupos (ligandos). Es útil para oxoácidos, oxoaniones y compuestos de coordinación: \ch{H2SO4} = \ch{[SO2(OH)2]} dihidroxidodioxidoazufre.
    \item \negrita{Nombres vulgares admitidos}. Son nombres tradicionales que la IUPAC sigue aceptando: agua, amoniaco, metano, ácido sulfúrico, nitrato\dots \textbf{Solo} se admiten los que están en la lista oficial (ver apartado~\ref{sec:vulgares}); la antigua nomenclatura tradicional (\mal{anhídrido sulfúrico}, \mal{óxido férrico}\dots) ya \negrita{no} es válida.
\end{enumerate}

%=====================================================================
\section{Compuestos binarios}\label{sec:binarios}
%=====================================================================

Los compuestos binarios están formados por \negrita{dos elementos}. Todos se nombran con la misma receta.

\subsection{La receta general}

\esquemafig{esq-binario}

El elemento de la derecha se nombra con la \negrita{raíz de su nombre + uro} (si es el oxígeno, se dice \negrita{óxido}). En el apartado siguiente se explica cómo obtener la raíz.

\figura{fig-receta}

\subsection{Cómo se obtiene la raíz del nombre de un elemento}\label{sec:raiz}

Para nombrar el elemento de la derecha se usa la \negrita{raíz} de su nombre seguida de \blueremark{-uro}. La raíz es lo que queda al \negrita{quitar la terminación} del nombre del elemento (\blueremark{-o}, \blueremark{-io}, \blueremark{-ico}, \blueremark{-ono}, \blueremark{-ógeno}, \blueremark{-oro}\dots).

\figura{fig-raiz}

\begin{center}
    \SetTblrInner{rowsep=3pt}
    % --- tabla antigua ---
    % \begin{longtblr}[caption={Raíz del nombre de los no metales más usados}]
        % {colspec={Q[l,m] Q[c,m] Q[c,m] Q[l,m]}, row{1}={bg=myOrange3!90!white, fg=white, font=\sffamily}, rowhead=1}
        % \hline
        % Elemento & Se quita & Raíz & Nombre con \blueremark{-uro} \\
        % \hline
        % flúor & (nada) & fluor & \negrita{fluor}uro \\
        % cloro & -o & clor & \negrita{clor}uro \\
        % bromo & -o & brom & \negrita{brom}uro \\
        % yodo & -o & yod & \negrita{yod}uro \\
        % boro & -o & bor & \negrita{bor}uro \\
        % selenio & -o & seleni & \negrita{seleni}uro \\
        % telurio & -io & telur & \negrita{telur}uro \\
        % nitrógeno & -ógeno & nitr & \negrita{nitr}uro \\
        % hidrógeno & -ógeno & hidr & \negrita{hidr}uro \\
        % carbono & -ono & carb & \negrita{carb}uro \\
        % fósforo & -oro & fosf & \negrita{fosf}uro \\
        % arsénico & -ico & arsen & \negrita{arsen}uro \\
        % \SetRow{myOrange3!15!white} azufre & raíz latina (\emph{sulphur}) & sulf & \negrita{sulf}uro \\
        % \SetRow{myOrange3!15!white} oxígeno & excepción & --- & \negrita{óxido} \\
        % \hline
    % \end{longtblr}
    % --- tabla nueva ---
    \begin{longtable}{@{}lccl@{}}
        \caption{Raíz del nombre de los no metales más usados} \\
        \toprule
        \textbf{Elemento} & \textbf{Se quita} & \textbf{Raíz} & \textbf{Nombre con \blueremark{-uro}} \\
        \midrule
        \endfirsthead
        \toprule
        \textbf{Elemento} & \textbf{Se quita} & \textbf{Raíz} & \textbf{Nombre con \blueremark{-uro}} \\
        \midrule
        \endhead
        \bottomrule
        \endlastfoot
        flúor & (nada) & fluor & \negrita{fluor}uro \\[3pt]
        cloro & -o & clor & \negrita{clor}uro \\[3pt]
        bromo & -o & brom & \negrita{brom}uro \\[3pt]
        yodo & -o & yod & \negrita{yod}uro \\[3pt]
        boro & -o & bor & \negrita{bor}uro \\[3pt]
        selenio & -o & seleni & \negrita{seleni}uro \\[3pt]
        telurio & -io & telur & \negrita{telur}uro \\[3pt]
        nitrógeno & -ógeno & nitr & \negrita{nitr}uro \\[3pt]
        hidrógeno & -ógeno & hidr & \negrita{hidr}uro \\[3pt]
        carbono & -ono & carb & \negrita{carb}uro \\[3pt]
        fósforo & -oro & fosf & \negrita{fosf}uro \\[3pt]
        arsénico & -ico & arsen & \negrita{arsen}uro \\[3pt]
        azufre & raíz latina (\emph{sulphur}) & sulf & \negrita{sulf}uro \\[3pt]
        oxígeno & excepción & --- & \negrita{óxido} \\
    \end{longtable}
\end{center}

\begin{consejos}{la raíz del nombre}
    \begin{enumerate}[label=\numberballgreen]
        \item Escribe el nombre del elemento y \negrita{tacha la terminación}: clor\mal{o}, nitr\mal{ógeno}, carb\mal{ono}.
        \item Añade \negrita{-uro} a lo que queda: clor + uro = \negrita{cloruro}.
        \item La raíz pierde la tilde: fósforo \(\to\) \negrita{fosf}uro; arsénico \(\to\) \negrita{arsen}uro.
        \item Dos casos especiales que hay que memorizar: el azufre usa su raíz latina (\negrita{sulf}uro) y el oxígeno no lleva -uro: se dice \negrita{óxido}.
    \end{enumerate}
\end{consejos}

\begin{exercise}
    Escribe la raíz y el nombre terminado en -uro (u óxido) de: \textbf{a)} cloro; \textbf{b)} yodo; \textbf{c)} nitrógeno; \textbf{d)} carbono; \textbf{e)} fósforo; \textbf{f)} azufre; \textbf{g)} hidrógeno; \textbf{h)} oxígeno.
    \begin{solu}
        \textbf{a)} clor \(\to\) \textcolor{red}{cloruro}; \textbf{b)} yod \(\to\) \textcolor{red}{yoduro}; \textbf{c)} nitr \(\to\) \textcolor{red}{nitruro}; \textbf{d)} carb \(\to\) \textcolor{red}{carburo}; \textbf{e)} fosf \(\to\) \textcolor{red}{fosfuro}; \textbf{f)} sulf \(\to\) \textcolor{red}{sulfuro}; \textbf{g)} hidr \(\to\) \textcolor{red}{hidruro}; \textbf{h)} \textcolor{red}{óxido} (excepción)
    \end{solu}
\end{exercise}

\subsection{Consejos para nombrar y formular}

\begin{consejos}{compuestos binarios}
    \begin{enumerate}[label=\numberballgreen]
        \item \negrita{De la fórmula al nombre}: lee la fórmula de \negrita{derecha a izquierda} y traduce cada subíndice a un prefijo.
        \item \negrita{Del nombre a la fórmula}: el nombre ya te dice los subíndices. \negrita{di}fosfuro de \negrita{tri}calcio \(\to\) \ch{Ca3P2}. Escribe a la izquierda lo que va después de «de».
        \item \negrita{Sin prefijos} (óxido de aluminio, cloruro de cinc\dots): usa los números de oxidación de la figura del apartado~\ref{sec:eo} y la regla de la suma cero.
    \end{enumerate}
\end{consejos}

\subsection{Óxidos}\label{sec:oxidos}

Son compuestos formados por \negrita{oxígeno} y otro elemento (metal o no metal). El oxígeno actúa con número de oxidación \azul{\(-2\)} y se escribe a la derecha. Los compuestos del oxígeno con los halógenos \negrita{no} son óxidos (apartado~\ref{sec:halogenuros_oxigeno}).

\formulan
\begin{enumerate}[label=\itembolascolor{\arabic*}{green!50!black}]
    \item Escribimos el elemento a la izquierda y el oxígeno a la derecha.
    \item Si el nombre lleva prefijos, los convertimos directamente en subíndices.
    \item Si el nombre lleva número de oxidación o no lleva nada, buscamos el número mínimo de átomos que hace que la suma sea cero.
\end{enumerate}

\ejemplos
\ballnumber{1} : Óxido de sodio: \ox{+1,Na} \ox{-2,O} \(\to\) \(2\cdot(+1) + (-2) = 0\) \(\to\) \ch{Na2O}

\ballnumber{2} : Óxido de calcio: \ox{+2,Ca} \ox{-2,O} \(\to\) \((+2) + (-2) = 0\) \(\to\) \ch{CaO}

\ballnumber{3} : Óxido de hierro(III): \ox{+3,Fe} \ox{-2,O} \(\to\) \(2\cdot(+3) + 3\cdot(-2) = 0\) \(\to\) \ch{Fe2O3}

\ballnumber{4} : Pentaóxido de difósforo \(\to\) \ch{P2O5}

\nombran
\begin{enumerate}[label=\itembolascolor{\arabic*}{blue}]
    \item \textbf{\uline{Con prefijos multiplicadores}} (forma preferida):
          \esquemafig{esq-oxido-pref}
          \begin{enumerate}[label=\textcolor{red}{\(\bullet\)}]
              \item Prefijo\subin{1}: número de átomos de oxígeno.
              \item Prefijo\subin{2}: número de átomos del otro elemento.
          \end{enumerate}

          \ejemplos
          \begin{center}
              \SetTblrInner{rowsep=3pt}
              % --- tabla antigua ---
              % \begin{tblr}{colspec={Q[c,m] Q[l,m]}, row{1}={bg=myOrange3!90!white, fg=white, font=\sffamily}, hline{1,2,Z}={1pt}, hline{3-Y}={dashed}}
                  % Fórmula & Nombre \\
                  % \ch{CO} & monóxido de carbono \\
                  % \ch{CO2} & dióxido de carbono \\
                  % \ch{Na2O} & óxido de disodio \\
                  % \ch{Fe2O3} & trióxido de dihierro \\
                  % \ch{Ni2O3} & trióxido de diníquel \\
                  % \ch{SO3} & trióxido de azufre \\
                  % \ch{N2O5} & pentaóxido de dinitrógeno \\
                  % \ch{CaO} & óxido de calcio \\
                  % \ch{P4O10} & decaóxido de tetrafósforo \\
              % \end{tblr}
              % --- tabla nueva ---
              \begin{tabular}{@{}cl@{}}
                  \toprule
                  \textbf{Fórmula} & \textbf{Nombre} \\
                  \midrule
                  \ch{CO} & monóxido de carbono \\[3pt]
                  \ch{CO2} & dióxido de carbono \\[3pt]
                  \ch{Na2O} & óxido de disodio \\[3pt]
                  \ch{Fe2O3} & trióxido de dihierro \\[3pt]
                  \ch{Ni2O3} & trióxido de diníquel \\[3pt]
                  \ch{SO3} & trióxido de azufre \\[3pt]
                  \ch{N2O5} & pentaóxido de dinitrógeno \\[3pt]
                  \ch{CaO} & óxido de calcio \\[3pt]
                  \ch{P4O10} & decaóxido de tetrafósforo \\
                  \bottomrule
              \end{tabular}
          \end{center}

    \item \textbf{\uline{Con el número de oxidación}} \conox:
          \esquemafig{esq-oxido-eo}
          El número de oxidación va pegado al nombre, sin espacio: hierro(III). Si el elemento tiene un único número de oxidación (alcalinos, alcalinotérreos, Zn, Ag, Al) se omite.

          \ejemplos
          \ballnumber{1} : \ch{Ni2O3}: óxido de níquel(III)

          \ballnumber{2} : \ch{SO3}: óxido de azufre(VI)

          \ballnumber{3} : \ch{Na2O}: \mal{óxido de sodio(I)} \(\to\) \textcolor{blue}{óxido de sodio}, porque el sodio solo actúa con \(+1\).

    \item \textbf{\uline{Con el número de carga}} \concarga{} \ampliacion: solo para los óxidos de metales, que son iónicos. \ch{Fe2O3}: óxido de hierro(3+); \ch{Cu2O}: óxido de cobre(1+). Los óxidos de no metales (\ch{SO3}, \ch{CO2}\dots) son covalentes y \negrita{no} se pueden nombrar así.
\end{enumerate}

\begin{marker}
    \negrita{No} se debe usar la nomenclatura antigua: \mal{anhídrido sulfúrico}, \mal{óxido férrico}, \mal{óxido ferroso}. Para el \ch{SO3} se dice trióxido de azufre u óxido de azufre(VI).
\end{marker}

\begin{exercise}
    \textbf{1.} Nombra cada óxido de dos formas, \conpref{} y \conox: \textbf{a)} \ch{CO}; \textbf{b)} \ch{SO2}; \textbf{c)} \ch{N2O3}; \textbf{d)} \ch{P2O5}; \textbf{e)} \ch{PbO2}; \textbf{f)} \ch{CaO}; \textbf{g)} \ch{Cu2O}.\\
    \textbf{2.} \ampliacion{} Nombra con \concarga{} los óxidos del apartado anterior que lo admitan. ¿Por qué los demás no lo admiten?
    \begin{solu}
        \textbf{1. a)} \conpref{} \textcolor{red}{monóxido de carbono} \quad \conox{} \textcolor{red}{óxido de carbono(II)}\\
        \textbf{b)} \conpref{} \textcolor{red}{dióxido de azufre} \quad \conox{} \textcolor{red}{óxido de azufre(IV)}\\
        \textbf{c)} \conpref{} \textcolor{red}{trióxido de dinitrógeno} \quad \conox{} \textcolor{red}{óxido de nitrógeno(III)}\\
        \textbf{d)} \conpref{} \textcolor{red}{pentaóxido de difósforo} \quad \conox{} \textcolor{red}{óxido de fósforo(V)}\\
        \textbf{e)} \conpref{} \textcolor{red}{dióxido de plomo} \quad \conox{} \textcolor{red}{óxido de plomo(IV)}\\
        \textbf{f)} \conpref{} \textcolor{red}{óxido de calcio} \quad \conox{} \textcolor{red}{óxido de calcio} (se omite porque el calcio solo actúa con \(+2\))\\
        \textbf{g)} \conpref{} \textcolor{red}{óxido de dicobre} \quad \conox{} \textcolor{red}{óxido de cobre(I)}\\
        \textbf{2.} \concarga{} \textbf{e)} \textcolor{red}{óxido de plomo(4+)}; \textbf{f)} \textcolor{red}{óxido de calcio}; \textbf{g)} \textcolor{red}{óxido de cobre(1+)}.\\
        \textcolor{blue}{\ch{CO}, \ch{SO2}, \ch{N2O3} y \ch{P2O5} son óxidos de no metales: son covalentes, no tienen iones y por tanto no tienen número de carga.}
    \end{solu}
\end{exercise}

\begin{exercise}
    Formula: \textbf{a)} dióxido de nitrógeno; \textbf{b)} óxido de aluminio; \textbf{c)} pentaóxido de difósforo; \textbf{d)} óxido de hierro(II); \textbf{e)} dióxido de silicio; \textbf{f)} trióxido de azufre; \textbf{g)} óxido de cobre(2+).
    \begin{solu}
        \textbf{a)} \textcolor{red}{\ch{NO2}}\\
        \textbf{b)} \textcolor{blue}{Al \(+3\), O \(-2\): \(2\cdot(+3)+3\cdot(-2)=0\)} \(\to\) \textcolor{red}{\ch{Al2O3}}\\
        \textbf{c)} \textcolor{red}{\ch{P2O5}}\\
        \textbf{d)} \conox{} \textcolor{blue}{Fe \(+2\), O \(-2\)} \(\to\) \textcolor{red}{\ch{FeO}}\\
        \textbf{e)} \textcolor{red}{\ch{SiO2}}\\
        \textbf{f)} \textcolor{red}{\ch{SO3}}\\
        \textbf{g)} \concarga{} \textcolor{blue}{iones \ch{Cu^2+} y \ch{O^2-}} \(\to\) \textcolor{red}{\ch{CuO}}
    \end{solu}
\end{exercise}

\subsection{Oxígeno con halógenos: halogenuros de oxígeno}\label{sec:halogenuros_oxigeno}

En la secuencia de los elementos los halógenos (F, Cl, Br, I) están a la \negrita{derecha} del oxígeno. Por eso, en sus compuestos con oxígeno, el oxígeno se escribe a la \negrita{izquierda} y el halógeno se nombra con \negrita{raíz + uro}: \negrita{no son óxidos}, sino \negrita{halogenuros de oxígeno}.

\esquemafig{esq-halogenuro-oxigeno}

\ejemplos
\begin{center}
    \SetTblrInner{rowsep=3pt}
    % --- tabla antigua ---
    % \begin{tblr}{colspec={Q[c,m] Q[l,m]}, row{1}={bg=myOrange3!90!white, fg=white, font=\sffamily}, hline{1,2,Z}={1pt}, hline{3-Y}={dashed}}
        % Fórmula & Nombre \\
        % \ch{OF2} & difluoruro de oxígeno \\
        % \ch{OCl2} & dicloruro de oxígeno \\
        % \ch{O3Cl2} & dicloruro de trioxígeno \\
        % \ch{O5Br2} & dibromuro de pentaoxígeno \\
        % \ch{O7I2} & diyoduro de heptaoxígeno \\
    % \end{tblr}
    % --- tabla nueva ---
    \begin{tabular}{@{}cl@{}}
        \toprule
        \textbf{Fórmula} & \textbf{Nombre} \\
        \midrule
        \ch{OF2} & difluoruro de oxígeno \\[3pt]
        \ch{OCl2} & dicloruro de oxígeno \\[3pt]
        \ch{O3Cl2} & dicloruro de trioxígeno \\[3pt]
        \ch{O5Br2} & dibromuro de pentaoxígeno \\[3pt]
        \ch{O7I2} & diyoduro de heptaoxígeno \\
        \bottomrule
    \end{tabular}
\end{center}

\begin{marker}
    Es un error frecuente escribir \mal{\ch{Cl2O}} y llamarlo \mal{óxido de dicloro}. Según la IUPAC (2005) es \ch{OCl2}, \negrita{dicloruro de oxígeno}.
\end{marker}

\begin{exercise}
    \textbf{1.} Nombra con prefijos: \textbf{a)} \ch{OF2}; \textbf{b)} \ch{OBr2}; \textbf{c)} \ch{O3Cl2}; \textbf{d)} \ch{O7Cl2}.\\
    \textbf{2.} Formula: \textbf{a)} dicloruro de oxígeno; \textbf{b)} dibromuro de trioxígeno; \textbf{c)} diyoduro de pentaoxígeno.
    \begin{solu}
        \textcolor{blue}{El halógeno se nombra primero (raíz + uro) y el oxígeno se escribe a la izquierda.}\\
        \textbf{1. a)} \textcolor{red}{difluoruro de oxígeno}; \textbf{b)} \textcolor{red}{dibromuro de oxígeno}; \textbf{c)} \textcolor{red}{dicloruro de trioxígeno}; \textbf{d)} \textcolor{red}{dicloruro de heptaoxígeno}\\
        \textbf{2. a)} \textcolor{red}{\ch{OCl2}}; \textbf{b)} \textcolor{red}{\ch{O3Br2}}; \textbf{c)} \textcolor{red}{\ch{O5I2}}
    \end{solu}
\end{exercise}

\subsection[Peróxidos \ampliacionicono]{Peróxidos \ampliacion}\label{sec:peroxidos}

Están formados por un catión (normalmente un metal alcalino o alcalinotérreo, o el hidrógeno) y el \negrita{ion peróxido} \ch{O2^2-}, en el que los dos oxígenos están unidos entre sí. En los peróxidos el oxígeno actúa con número de oxidación \(-1\).

\begin{marker}
    ¿Óxido o peróxido? Haz la suma cero: si el oxígeno sale con \(-2\) es un óxido; si sale con \(-1\), es un peróxido. El grupo \ch{O2} \negrita{no se simplifica nunca}: \ch{Na2O2} y no \mal{\ch{NaO}}.
\end{marker}

\formulan
\begin{enumerate}[label=\itembolascolor{\arabic*}{green!40!black}]
    \item Escribimos el metal a la izquierda y el grupo peróxido \ch{O2^2-} a la derecha.
    \item Ponemos los átomos del metal necesarios para neutralizar la carga \(2-\) del grupo.
\end{enumerate}

\ejemplos
\ballnumber{1} : Peróxido de sodio: \ch{Na+} y \ch{O2^2-} \(\to\) \(2\cdot(+1) + (-2) = 0\) \(\to\) \ch{Na2O2}

\ballnumber{2} : Peróxido de bario: \ch{Ba^2+} y \ch{O2^2-} \(\to\) \ch{BaO2} (el 2 es del grupo peróxido, no se puede quitar).

\nombran
\begin{enumerate}[label=\itembolascolor{\arabic*}{blue}]
    \item \textbf{\uline{Con prefijos multiplicadores}}: igual que los óxidos. \ch{Na2O2}: dióxido de disodio; \ch{Li2O2}: dióxido de dilitio.
    \item \textbf{\uline{Nombre de peróxido}}:
          \esquemafig{esq-peroxido}
          \ch{Na2O2}: peróxido de sodio; \ch{CaO2}: peróxido de calcio; \ch{H2O2}: peróxido de hidrógeno (nombre vulgar admitido: \negrita{agua oxigenada}).
\end{enumerate}

\begin{exercise}
    \textbf{1.} Indica si son óxidos o peróxidos, justificándolo con la suma cero: \textbf{a)} \ch{K2O2}; \textbf{b)} \ch{BaO}; \textbf{c)} \ch{BaO2}; \textbf{d)} \ch{Li2O}.\\
    \textbf{2.} Formula: \textbf{a)} peróxido de litio; \textbf{b)} peróxido de calcio; \textbf{c)} peróxido de hidrógeno.
    \begin{solu}
        \textbf{1. a)} \(2\cdot(+1) + 2x = 0 \Rightarrow x = -1\): \textcolor{red}{peróxido} de potasio\\
        \textbf{b)} \((+2) + x = 0 \Rightarrow x = -2\): \textcolor{red}{óxido} de bario\\
        \textbf{c)} \((+2) + 2x = 0 \Rightarrow x = -1\): \textcolor{red}{peróxido} de bario\\
        \textbf{d)} \(2\cdot(+1) + x = 0 \Rightarrow x = -2\): \textcolor{red}{óxido} de litio\\
        \textbf{2. a)} \textcolor{red}{\ch{Li2O2}}; \textbf{b)} \textcolor{red}{\ch{CaO2}} (no se simplifica el grupo \ch{O2}); \textbf{c)} \textcolor{red}{\ch{H2O2}}
    \end{solu}
\end{exercise}

\subsection{Hidruros metálicos (metal + hidrógeno)}\label{sec:hidruros}

Son compuestos de un \negrita{metal} con hidrógeno. Con los metales, el hidrógeno actúa con número de oxidación \azul{\(-1\)} y se escribe a la derecha.

\formulan
\begin{enumerate}[label=\itembolascolor{\arabic*}{mygreen}]
    \item Ponemos el metal a la izquierda y el hidrógeno a la derecha.
    \item Como el hidrógeno es \(-1\), el número de hidrógenos coincide con el número de oxidación del metal.
\end{enumerate}

\ejemplos
\ballnumber{1} : Hidruro de sodio: \ox{+1,Na} \ox{-1,H} \(\to\) \ch{NaH}

\ballnumber{2} : Hidruro de calcio: \ox{+2,Ca} \ox{-1,H} \(\to\) \ch{CaH2}

\nombran
\begin{enumerate}[label=\itembolascolor{\arabic*}{blue}]
    \item \textbf{\uline{Con prefijos multiplicadores}}:
          \esquemafig{esq-hidruro-pref}
          El prefijo indica el número de hidrógenos. Como en la fórmula hay siempre un solo átomo de metal, el metal no lleva prefijo.

          \ejemplos
          \ballnumber{1} : \ch{AlH3}: trihidruro de aluminio \quad \ballnumber{2} : \ch{FeH2}: dihidruro de hierro \quad \ballnumber{3} : \ch{CoH3}: trihidruro de cobalto

    \item \textbf{\uline{Con el número de oxidación}} \conox:
          \esquemafig{esq-hidruro-eo}

          \ejemplos
          \ballnumber{1} : \ch{NiH2}: hidruro de níquel(II) \quad \ballnumber{2} : \ch{PbH4}: hidruro de plomo(IV) \quad \ballnumber{3} : \ch{ZnH2}: hidruro de cinc

    \item \textbf{\uline{Con el número de carga}} \concarga{} \ampliacion: \ch{NiH2}: hidruro de níquel(2+); \ch{FeH3}: hidruro de hierro(3+).
\end{enumerate}

\begin{exercise}
    \textbf{1.} Nombra cada hidruro de dos formas, \conpref{} y \conox: \textbf{a)} \ch{LiH}; \textbf{b)} \ch{CaH2}; \textbf{c)} \ch{FeH3}; \textbf{d)} \ch{PbH4}; \textbf{e)} \ch{CuH}.\\
    \textbf{2.} Formula: \textbf{a)} dihidruro de magnesio; \textbf{b)} hidruro de aluminio; \textbf{c)} hidruro de estaño(IV); \textbf{d)} trihidruro de cobalto; \textbf{e)} hidruro de níquel(3+).\\
    \textbf{3.} \ampliacion{} Nombra con \concarga: \textbf{a)} \ch{FeH3}; \textbf{b)} \ch{CuH}.
    \begin{solu}
        \textbf{1. a)} \conpref{} \textcolor{red}{hidruro de litio} \quad \conox{} \textcolor{red}{hidruro de litio} (se omite: el litio solo actúa con \(+1\))\\
        \textbf{b)} \conpref{} \textcolor{red}{dihidruro de calcio} \quad \conox{} \textcolor{red}{hidruro de calcio}\\
        \textbf{c)} \conpref{} \textcolor{red}{trihidruro de hierro} \quad \conox{} \textcolor{red}{hidruro de hierro(III)}\\
        \textbf{d)} \conpref{} \textcolor{red}{tetrahidruro de plomo} \quad \conox{} \textcolor{red}{hidruro de plomo(IV)}\\
        \textbf{e)} \conpref{} \textcolor{red}{hidruro de cobre} \quad \conox{} \textcolor{red}{hidruro de cobre(I)}\\
        \textbf{2.} \textcolor{blue}{Nº de H = nº de oxidación (o carga) del metal:} \textbf{a)} \textcolor{red}{\ch{MgH2}}; \textbf{b)} \textcolor{red}{\ch{AlH3}}; \textbf{c)} \textcolor{red}{\ch{SnH4}}; \textbf{d)} \textcolor{red}{\ch{CoH3}}; \textbf{e)} \textcolor{red}{\ch{NiH3}}\\
        \textbf{3.} \concarga{} \textbf{a)} \textcolor{red}{hidruro de hierro(3+)}; \textbf{b)} \textcolor{red}{hidruro de cobre(1+)}
    \end{solu}
\end{exercise}

\subsection{Hidrógeno con no metales}\label{sec:hidruros_nm}

Con los no metales el hidrógeno actúa con número de oxidación \azul{\(+1\)}. Hay que distinguir dos casos según el grupo del no metal.

\figura{fig-hidruros}

\subsubsection{Con los grupos 16 y 17: halogenuros y calcogenuros de hidrógeno}\label{sec:hidruros_16_17}

El hidrógeno se escribe a la \negrita{izquierda} (porque en la secuencia está antes que S, Se, Te, F, Cl, Br, I). El no metal actúa con su número de oxidación negativo: \(-2\) los del grupo 16 y \(-1\) los del grupo 17.

\esquemafig{esq-hidrogeno-nm}

\ejemplos
\ballnumber{1} : \ch{HF}: fluoruro de hidrógeno \quad \ballnumber{2} : \ch{HCl}: cloruro de hidrógeno \quad \ballnumber{3} : \ch{H2S}: sulfuro de dihidrógeno \quad \ballnumber{4} : \ch{H2Se}: seleniuro de dihidrógeno

Como no existe ambigüedad, también se admite omitir el prefijo del hidrógeno: \ch{H2S} sulfuro de hidrógeno.

\begin{marker}
    \negrita{¿Y el «ácido clorhídrico»?} Cuando estas sustancias se disuelven en agua, a la disolución se le ha llamado tradicionalmente así:
    \esquemafig{esq-hidracido}
    \ch{HCl_{(aq)}} «ácido clorhídrico», \ch{H2S_{(aq)}} «ácido sulfhídrico»\dots

    ¡Atención! Estos nombres corresponden a \negrita{disoluciones} (mezclas de composición variable), no a compuestos químicos, y por eso \negrita{la IUPAC no los admite}. El compuesto \ch{HCl} se llama siempre \negrita{cloruro de hidrógeno}.
\end{marker}

\begin{exercise}
    \textbf{1.} Nombra: \textbf{a)} \ch{HBr}; \textbf{b)} \ch{H2Se}; \textbf{c)} \ch{HI}; \textbf{d)} \ch{H2Te}; \textbf{e)} \ch{HF}.\\
    \textbf{2.} Formula: \textbf{a)} yoduro de hidrógeno; \textbf{b)} bromuro de hidrógeno; \textbf{c)} sulfuro de dihidrógeno.\\
    \textbf{3.} ¿Por qué no es correcto decir que el compuesto \ch{HCl} se llama ácido clorhídrico?
    \begin{solu}
        \textbf{1. a)} \textcolor{red}{bromuro de hidrógeno}; \textbf{b)} \textcolor{red}{seleniuro de dihidrógeno}; \textbf{c)} \textcolor{red}{yoduro de hidrógeno}; \textbf{d)} \textcolor{red}{telururo de dihidrógeno}; \textbf{e)} \textcolor{red}{fluoruro de hidrógeno}\\
        \textbf{2. a)} \textcolor{red}{\ch{HI}}; \textbf{b)} \textcolor{red}{\ch{HBr}}; \textbf{c)} \textcolor{red}{\ch{H2S}}\\
        \textbf{3.} \textcolor{blue}{Porque «ácido clorhídrico» es el nombre tradicional de la disolución de \ch{HCl} en agua, que es una mezcla de composición variable; la IUPAC no lo admite.} El compuesto \ch{HCl} se llama \textcolor{red}{cloruro de hidrógeno}.
    \end{solu}
\end{exercise}

\subsubsection{Con los grupos 13, 14 y 15: hidruros progenitores}

Con B, C, Si, N, P, As, Sb\dots el hidrógeno se escribe a la \negrita{derecha}. Estos compuestos se nombran preferentemente con su \negrita{nombre de hidruro progenitor} (nomenclatura de sustitución), que termina en \blueremark{-ano}; también pueden nombrarse con prefijos multiplicadores.

% --- tabla antigua ---
% \begin{longtblr}[caption={Hidruros de los no metales},
        % note{a} = {Nomenclatura de composición con prefijos multiplicadores},
        % note{b} = {Nombre del hidruro progenitor (nomenclatura de sustitución) o nombre vulgar admitido}
    % ]
    % {colspec={Q[c,m] Q[c,m] Q[c,m]},
        % row{1} = {bg=myOrange3!90!white, fg = white, font=\sffamily}
    % }
    % \hline
    % Compuesto & Prefijos multiplicadores \TblrNote{a} & Hidruro progenitor / vulgar \TblrNote{b} \\
    % \hline
    % \SetRow{myOrange3!20!white}\SetCell[c=3]{c} Los que hay que conocer\\
    % \hline
    % \ch{CH4} & tetrahidruro de carbono & \negrita{metano} \\
    % \ch{SiH4} & tetrahidruro de silicio & silano \\
    % \ch{NH3} & trihidruro de nitrógeno & \negrita{amoniaco}, azano \\
    % \ch{PH3} & trihidruro de fósforo & fosfano (\mal{fosfina}) \\
    % \ch{AsH3} & trihidruro de arsénico & arsano (\mal{arsina}) \\
    % \ch{SbH3} & trihidruro de antimonio & estibano (\mal{estibina}) \\
    % \ch{H2O} & --- & \negrita{agua}, oxidano \\
    % \hline
    % \SetRow{myOrange3!20!white} \SetCell[c=3]{c} \textcolor{myOrange3}{\faRocket} Ampliación\\
    % \hline
    % \ch{BH3} & trihidruro de boro & borano \\
    % \ch{AlH3} & trihidruro de aluminio & alumano \\
    % \ch{GaH3} & trihidruro de galio & galano \\
    % \ch{InH3} & trihidruro de indio & indigano \\
    % \ch{TlH3} & trihidruro de talio & talano \\
    % \ch{GeH4} & tetrahidruro de germanio & germano \\
    % \ch{SnH4} & tetrahidruro de estaño & estannano \\
    % \ch{PbH4} & tetrahidruro de plomo & plumbano\\
    % \ch{BiH3} & trihidruro de bismuto & bismutano\\
    % \ch{H2S} & sulfuro de dihidrógeno & sulfano\\
    % \ch{H2Se} & seleniuro de dihidrógeno & selano\\
    % \ch{H2Te} & telururo de dihidrógeno & telano\\
    % \ch{H2Po} & polonuro de dihidrógeno & polano\\
    % \ch{HF} & fluoruro de hidrógeno & fluorano\\
    % \ch{HCl} & cloruro de hidrógeno & clorano\\
    % \ch{HBr} & bromuro de hidrógeno & bromano\\
    % \ch{HI} & yoduro de hidrógeno & yodano\\
    % \ch{HAt} & astaturo de hidrógeno & astatano\\
    % \hline
% \end{longtblr}
% --- tabla nueva ---
\begin{longtable}{@{}ccc@{}}
    \caption{Hidruros de los no metales} \\
    \toprule
    \textbf{Compuesto} & \textbf{Prefijos multiplicadores \textsuperscript{a}} & \textbf{Hidruro progenitor / vulgar \textsuperscript{b}} \\
    \midrule
    \endfirsthead
    \toprule
    \textbf{Compuesto} & \textbf{Prefijos multiplicadores \textsuperscript{a}} & \textbf{Hidruro progenitor / vulgar \textsuperscript{b}} \\
    \midrule
    \endhead
    \bottomrule
    \endlastfoot
    \multicolumn{3}{c}{Los que hay que conocer} \\[3pt]
    \ch{CH4} & tetrahidruro de carbono & \negrita{metano} \\[3pt]
    \ch{SiH4} & tetrahidruro de silicio & silano \\[3pt]
    \ch{NH3} & trihidruro de nitrógeno & \negrita{amoniaco}, azano \\[3pt]
    \ch{PH3} & trihidruro de fósforo & fosfano (\mal{fosfina}) \\[3pt]
    \ch{AsH3} & trihidruro de arsénico & arsano (\mal{arsina}) \\[3pt]
    \ch{SbH3} & trihidruro de antimonio & estibano (\mal{estibina}) \\[3pt]
    \ch{H2O} & --- & \negrita{agua}, oxidano \\[3pt]
    \multicolumn{3}{c}{\textcolor{myOrange3}{\faRocket} Ampliación} \\[3pt]
    \ch{BH3} & trihidruro de boro & borano \\[3pt]
    \ch{AlH3} & trihidruro de aluminio & alumano \\[3pt]
    \ch{GaH3} & trihidruro de galio & galano \\[3pt]
    \ch{InH3} & trihidruro de indio & indigano \\[3pt]
    \ch{TlH3} & trihidruro de talio & talano \\[3pt]
    \ch{GeH4} & tetrahidruro de germanio & germano \\[3pt]
    \ch{SnH4} & tetrahidruro de estaño & estannano \\[3pt]
    \ch{PbH4} & tetrahidruro de plomo & plumbano \\[3pt]
    \ch{BiH3} & trihidruro de bismuto & bismutano \\[3pt]
    \ch{H2S} & sulfuro de dihidrógeno & sulfano \\[3pt]
    \ch{H2Se} & seleniuro de dihidrógeno & selano \\[3pt]
    \ch{H2Te} & telururo de dihidrógeno & telano \\[3pt]
    \ch{H2Po} & polonuro de dihidrógeno & polano \\[3pt]
    \ch{HF} & fluoruro de hidrógeno & fluorano \\[3pt]
    \ch{HCl} & cloruro de hidrógeno & clorano \\[3pt]
    \ch{HBr} & bromuro de hidrógeno & bromano \\[3pt]
    \ch{HI} & yoduro de hidrógeno & yodano \\[3pt]
    \ch{HAt} & astaturo de hidrógeno & astatano \\
\end{longtable}
{\footnotesize
    \textsuperscript{a}\,Nomenclatura de composición con prefijos multiplicadores\par
    \textsuperscript{b}\,Nombre del hidruro progenitor (nomenclatura de sustitución) o nombre vulgar admitido\par
}

\begin{exercise}
    \textbf{1.} Escribe el nombre preferido de: \textbf{a)} \ch{CH4}; \textbf{b)} \ch{NH3}; \textbf{c)} \ch{PH3}; \textbf{d)} \ch{SiH4}; \textbf{e)} \ch{AsH3}.\\
    \textbf{2.} Formula: \textbf{a)} estibano; \textbf{b)} borano; \textbf{c)} azano.
    \begin{solu}
        \textbf{1. a)} \textcolor{red}{metano}; \textbf{b)} \textcolor{red}{amoniaco} o azano; \textbf{c)} \textcolor{red}{fosfano} (no fosfina); \textbf{d)} \textcolor{red}{silano}; \textbf{e)} \textcolor{red}{arsano} (no arsina)\\
        \textbf{2. a)} \textcolor{red}{\ch{SbH3}}; \textbf{b)} \textcolor{red}{\ch{BH3}}; \textbf{c)} \textcolor{red}{\ch{NH3}}
    \end{solu}
\end{exercise}

\subsection{Sales binarias: metal + no metal}\label{sec:sales_binarias}

Son compuestos formados por un metal (a la izquierda) y un no metal distinto del oxígeno y del hidrógeno (a la derecha). El no metal actúa con su número de oxidación negativo (ver la figura del apartado~\ref{sec:eo}): F, Cl, Br, I \(\to -1\); S, Se, Te \(\to -2\); N, P \(\to -3\); C \(\to -4\).

\formulan
\begin{enumerate}[label=\itembolascolor{\arabic*}{green!40!black}]
    \item Ponemos el metal a la izquierda y el no metal a la derecha.
    \item Con prefijos: los prefijos son los subíndices.
    Sin prefijos o con número de oxidación: aplicamos la suma cero.
\end{enumerate}

\ejemplos
\ballnumber{1} : Cloruro de sodio: \ox{+1,Na} \ox{-1,Cl} \(\to\) \ch{NaCl}

\ballnumber{2} : Sulfuro de estaño(II): \ox{+2,Sn} \ox{-2,S} \(\to\) \((+2)+(-2)=0\) \(\to\) \ch{SnS}

\ballnumber{3} : Cloruro de cinc: \ox{+2,Zn} \ox{-1,Cl} \(\to\) \((+2)+2\cdot(-1)=0\) \(\to\) \ch{ZnCl2}

\nombran
\begin{enumerate}[label=\itembolascolor{\arabic*}{blue}]
    \item \textbf{\uline{Con prefijos multiplicadores}}:
          \esquemafig{esq-sal-pref}

          \ejemplos
          \ballnumber{1} : \ch{Al4C3}: tricarburo de tetraaluminio \quad \ballnumber{2} : \ch{SnTe2}: ditelururo de estaño

          \ballnumber{3} : \ch{Ca3P2}: difosfuro de tricalcio \quad \ballnumber{4} : \ch{AgBr}: bromuro de plata

    \item \textbf{\uline{Con el número de oxidación}} \conox:
          \esquemafig{esq-sal-eo}

          \ejemplos
          \ballnumber{1} : \ch{Al4C3}: carburo de aluminio \quad \ballnumber{2} : \ch{SnTe2}: telururo de estaño(IV) \quad \ballnumber{3} : \ch{FeCl3}: cloruro de hierro(III)

    \item \textbf{\uline{Con el número de carga}} \concarga{} \ampliacion{} (solo compuestos iónicos): \ch{FeCl2}: cloruro de hierro(2+); \ch{CuS}: sulfuro de cobre(2+).
\end{enumerate}

\begin{exercise}
    \textbf{1.} Nombra cada sal de dos formas, \conpref{} y \conox: \textbf{a)} \ch{NaBr}; \textbf{b)} \ch{CaF2}; \textbf{c)} \ch{FeS}; \textbf{d)} \ch{Al2S3}; \textbf{e)} \ch{CuCl2}; \textbf{f)} \ch{Mg3N2}; \textbf{g)} \ch{AgI}.\\
    \textbf{2.} Formula: \textbf{a)} cloruro de cinc; \textbf{b)} sulfuro de dipotasio; \textbf{c)} yoduro de plomo(II); \textbf{d)} nitruro de litio; \textbf{e)} tricloruro de hierro; \textbf{f)} bromuro de cobre(1+).\\
    \textbf{3.} \ampliacion{} Nombra con \concarga: \textbf{a)} \ch{FeS}; \textbf{b)} \ch{CuCl2}.
    \begin{solu}
        \textbf{1. a)} \conpref{} \textcolor{red}{bromuro de sodio} \quad \conox{} \textcolor{red}{bromuro de sodio}\\
        \textbf{b)} \conpref{} \textcolor{red}{difluoruro de calcio} \quad \conox{} \textcolor{red}{fluoruro de calcio}\\
        \textbf{c)} \conpref{} \textcolor{red}{monosulfuro de hierro} \quad \conox{} \textcolor{red}{sulfuro de hierro(II)}\\
        \textbf{d)} \conpref{} \textcolor{red}{trisulfuro de dialuminio} \quad \conox{} \textcolor{red}{sulfuro de aluminio}\\
        \textbf{e)} \conpref{} \textcolor{red}{dicloruro de cobre} \quad \conox{} \textcolor{red}{cloruro de cobre(II)}\\
        \textbf{f)} \conpref{} \textcolor{red}{dinitruro de trimagnesio} \quad \conox{} \textcolor{red}{nitruro de magnesio}\\
        \textbf{g)} \conpref{} \textcolor{red}{yoduro de plata} \quad \conox{} \textcolor{red}{yoduro de plata}\\
        \textcolor{blue}{Con \conox{} el número se omite cuando el metal solo tiene uno (Na, Ca, Al, Mg, Ag).}\\
        \textbf{2. a)} \textcolor{blue}{Zn \(+2\), Cl \(-1\)} \(\to\) \textcolor{red}{\ch{ZnCl2}}; \textbf{b)} \textcolor{red}{\ch{K2S}}; \textbf{c)} \textcolor{red}{\ch{PbI2}}; \textbf{d)} \textcolor{blue}{Li \(+1\), N \(-3\)} \(\to\) \textcolor{red}{\ch{Li3N}}; \textbf{e)} \textcolor{red}{\ch{FeCl3}}; \textbf{f)} \textcolor{blue}{iones \ch{Cu+} y \ch{Br-}} \(\to\) \textcolor{red}{\ch{CuBr}}\\
        \textbf{3.} \concarga{} \textbf{a)} \textcolor{red}{sulfuro de hierro(2+)}; \textbf{b)} \textcolor{red}{cloruro de cobre(2+)}
    \end{solu}
\end{exercise}

\subsection{Sales binarias: no metal + no metal}\label{sec:nometal_nometal}

Son compuestos covalentes entre dos no metales (distintos del oxígeno y del hidrógeno). El que está más a la derecha en la secuencia de los elementos se escribe a la derecha y se nombra con la raíz de su nombre + \blueremark{-uro}. Se nombran \negrita{siempre con prefijos} (también puede usarse el número de oxidación, pero nunca el de carga, porque no son iónicos).

\ejemplos
\ballnumber{1} : \ch{PCl3}: tricloruro de fósforo \quad \ballnumber{2} : \ch{PCl5}: pentacloruro de fósforo \quad \ballnumber{3} : \ch{SF6}: hexafluoruro de azufre

\ballnumber{4} : \ch{SF4}: tetrafluoruro de azufre o fluoruro de azufre(IV) \quad \ballnumber{5} : \ch{CS2}: disulfuro de carbono \quad \ballnumber{6} : \ch{BrF3}: trifluoruro de bromo

\begin{exercise}
    \textbf{1.} Nombra \conpref: \textbf{a)} \ch{PBr3}; \textbf{b)} \ch{CCl4}; \textbf{c)} \ch{NCl3}; \textbf{d)} \ch{BF3}; \textbf{e)} \ch{IF7}.\\
    \textbf{2.} Formula: \textbf{a)} pentafluoruro de fósforo; \textbf{b)} disulfuro de carbono; \textbf{c)} tetrafluoruro de silicio.\\
    \textbf{3.} Nombra \conox{} el \ch{SF4} y el \ch{PCl3}. ¿Se podrían nombrar con \concarga?
    \begin{solu}
        \textbf{1. a)} \textcolor{red}{tribromuro de fósforo}; \textbf{b)} \textcolor{red}{tetracloruro de carbono}; \textbf{c)} \textcolor{red}{tricloruro de nitrógeno}; \textbf{d)} \textcolor{red}{trifluoruro de boro}; \textbf{e)} \textcolor{red}{heptafluoruro de yodo}\\
        \textbf{2. a)} \textcolor{red}{\ch{PF5}}; \textbf{b)} \textcolor{red}{\ch{CS2}}; \textbf{c)} \textcolor{red}{\ch{SiF4}}\\
        \textbf{3.} \conox{} \textcolor{red}{fluoruro de azufre(IV)}; \textcolor{red}{cloruro de fósforo(III)}. \textcolor{blue}{No se pueden nombrar con \concarga{} porque son compuestos covalentes: no tienen iones.}
    \end{solu}
\end{exercise}

%=====================================================================
\section{Hidróxidos}\label{sec:hidroxidos}
%=====================================================================

Los hidróxidos son compuestos ternarios (tres elementos), aunque se nombran como los binarios. Tienen carácter básico, por eso también se llaman bases (aunque no todas las bases son hidróxidos).

Están formados por un catión metálico y el \negrita{ion hidróxido} \ch{OH-}, que es un grupo indivisible con \negrita{número de carga} \(1-\). Su fórmula general es \ch{M(OH)_n}, donde \(n\) es el número de oxidación del metal.

\formulan
\begin{enumerate}[label=\itembolascolor{\arabic*}{mygreen}]
    \item Ponemos el metal a la izquierda y el grupo \ch{OH} a la derecha.
    \item Ponemos tantos grupos \ch{OH} como indique el número de oxidación del metal (suma cero).
    \item Si hay más de un grupo, \ch{OH} va entre paréntesis: \ch{Ca(OH)2}. Si solo hay uno, sin paréntesis: \ch{NaOH}.
\end{enumerate}

\ejemplos
\ballnumber{1}: Hidróxido de sodio: \ch{Na+} y \ch{OH-} \(\to\) \ch{NaOH}

\ballnumber{2}: Dihidróxido de calcio: \ch{Ca^2+} y \ch{OH-} \(\to\) \((+2)+2\cdot(-1)=0\) \(\to\) \ch{Ca(OH)2}

\nombran
\begin{enumerate}[label=\itembolascolor{\arabic*}{blue}]
    \item \textbf{\uline{Con prefijos multiplicadores}}:
          \esquemafig{esq-hidroxido-pref}
          El prefijo indica el número de grupos \ch{OH}.

          \ejemplos
          \ballnumber{1} : \ch{LiOH}: hidróxido de litio \quad \ballnumber{2} : \ch{Mg(OH)2}: dihidróxido de magnesio \quad \ballnumber{3} : \ch{Al(OH)3}: trihidróxido de aluminio

    \item \textbf{\uline{Con el número de oxidación}} \conox:
          \esquemafig{esq-hidroxido-eo}

          \ejemplos
          \ballnumber{1} : \ch{Fe(OH)2}: hidróxido de hierro(II) \quad \ballnumber{2} : \ch{CuOH}: hidróxido de cobre(I)

    \item \textbf{\uline{Con el número de carga}} \concarga{} \ampliacion: los hidróxidos son iónicos, así que también se puede indicar la carga del catión metálico.

          \ejemplos
          \ballnumber{1} : \ch{Fe(OH)2}: hidróxido de hierro(2+) \quad \ballnumber{2} : \ch{CuOH}: hidróxido de cobre(1+)
\end{enumerate}

\begin{exercise}
    \textbf{1.} Nombra cada hidróxido de dos formas, \conpref{} y \conox: \textbf{a)} \ch{KOH}; \textbf{b)} \ch{Ba(OH)2}; \textbf{c)} \ch{Fe(OH)3}; \textbf{d)} \ch{Pb(OH)4}; \textbf{e)} \ch{CuOH}.\\
    \textbf{2.} Formula: \textbf{a)} hidróxido de litio; \textbf{b)} dihidróxido de cinc; \textbf{c)} hidróxido de cobalto(II); \textbf{d)} hidróxido de níquel(3+).\\
    \textbf{3.} \ampliacion{} Nombra con \concarga: \textbf{a)} \ch{Fe(OH)3}; \textbf{b)} \ch{Pb(OH)4}; \textbf{c)} \ch{CuOH}.
    \begin{solu}
        \textbf{1. a)} \conpref{} \textcolor{red}{hidróxido de potasio} \quad \conox{} \textcolor{red}{hidróxido de potasio}\\
        \textbf{b)} \conpref{} \textcolor{red}{dihidróxido de bario} \quad \conox{} \textcolor{red}{hidróxido de bario}\\
        \textbf{c)} \conpref{} \textcolor{red}{trihidróxido de hierro} \quad \conox{} \textcolor{red}{hidróxido de hierro(III)}\\
        \textbf{d)} \conpref{} \textcolor{red}{tetrahidróxido de plomo} \quad \conox{} \textcolor{red}{hidróxido de plomo(IV)}\\
        \textbf{e)} \conpref{} \textcolor{red}{hidróxido de cobre} \quad \conox{} \textcolor{red}{hidróxido de cobre(I)}\\
        \textbf{2.} \textcolor{blue}{Tantos \ch{OH} como el número de oxidación (o la carga) del metal:} \textbf{a)} \textcolor{red}{\ch{LiOH}}; \textbf{b)} \textcolor{red}{\ch{Zn(OH)2}}; \textbf{c)} \conox{} \textcolor{red}{\ch{Co(OH)2}}; \textbf{d)} \concarga{} \textcolor{red}{\ch{Ni(OH)3}}\\
        \textbf{3.} \concarga{} \textbf{a)} \textcolor{red}{hidróxido de hierro(3+)}; \textbf{b)} \textcolor{red}{hidróxido de plomo(4+)}; \textbf{c)} \textcolor{red}{hidróxido de cobre(1+)}
    \end{solu}
\end{exercise}

%=====================================================================
\section{Iones}\label{sec:iones}
%=====================================================================

Los iones son átomos o grupos de átomos con carga eléctrica. En la fórmula la carga se escribe arriba a la derecha con el número delante del signo: \ch{Ca^2+}, \ch{S^2-}, \ch{Fe^3+} (y no \mal{\ch{Ca^{+2}}}).

\subsection{Cationes (iones positivos)}

Se forman cuando un átomo o grupo de átomos \negrita{pierde} electrones.

\subsubsection{Cationes monoatómicos}

Los forman sobre todo los \negrita{metales}. Se nombran con el nombre del elemento seguido del número de carga \concarga{} entre paréntesis, sin espacio (en los iones \negrita{no} se usa el número de oxidación):
\esquemafig{esq-cation}

El número de carga se escribe con el número delante del signo: (1+), (2+)\dots{} (números de Ewens--Bassett). Es incorrecto escribir \mal{(+1)}, \mal{(+2)}\dots

El número de carga puede omitirse cuando el elemento solo forma un catión (alcalinos, alcalinotérreos, Zn, Ag, Al): \ch{Na+} ion sodio.

\ejemplos
\noindent \ballnumber{1} : \ch{K+}: potasio(1+) o ion potasio\\
\ballnumber{2} : \ch{Ca^2+}: calcio(2+) o ion calcio\\
\ballnumber{3} : \ch{Fe^3+}: hierro(3+)\\
\ballnumber{4} : \ch{Mn^4+}: manganeso(4+)\\
\ballnumber{5} : \ch{H+}: hidrógeno(1+) o \blueremark{hidrón}\footnote{Tradicionalmente al ion \ch{H+} se le ha llamado protón (si al átomo de H le quitamos el \ch{e-} queda solo un protón). Sin embargo, la IUPAC reserva este nombre para el \ch{^1H+} y recomienda el término \textbf{hidrón} para el catión del hidrógeno natural.}

\subsubsection{Cationes poliatómicos: adición de hidrones (\ch{H+})}

Los hidruros progenitores de los grupos 15 y 16 pueden aceptar un ion \ch{H+}. Se nombran añadiendo la terminación \blueremark{-io} al nombre del hidruro progenitor (azano \(\to\) azan\negrita{io}).

\ejemplos
\noindent \ballnumber{1} : \ch{NH3 + H+ -> NH4+}: \negrita{amonio} (nombre vulgar admitido) o azanio\\
\ballnumber{2} : \ch{H2O + H+ -> H3O+}: \negrita{oxonio} (nombre vulgar admitido) u oxidanio\footnote{La IUPAC admite oxonio pero NO \mal{hidronio}.}\\
\ballnumber{3} : \ch{PH3 + H+ -> PH4+}: fosfanio

\begin{exercise}
    \textbf{1.} Nombra los cationes indicando su \concarga{} (los poliatómicos, con su nombre vulgar): \textbf{a)} \ch{Na+}; \textbf{b)} \ch{Mg^2+}; \textbf{c)} \ch{Fe^2+}; \textbf{d)} \ch{Cu+}; \textbf{e)} \ch{NH4+}; \textbf{f)} \ch{H3O+}.\\
    \textbf{2.} Escribe la fórmula: \textbf{a)} aluminio(3+); \textbf{b)} ion plomo(4+); \textbf{c)} fosfanio.
    \begin{solu}
        \textbf{1.} \concarga{} \textbf{a)} \textcolor{red}{sodio(1+)} o ion sodio; \textbf{b)} \textcolor{red}{magnesio(2+)} o ion magnesio; \textbf{c)} \textcolor{red}{hierro(2+)} (no se puede omitir: el hierro forma \ch{Fe^2+} y \ch{Fe^3+}); \textbf{d)} \textcolor{red}{cobre(1+)}; \textbf{e)} \textcolor{red}{amonio}; \textbf{f)} \textcolor{red}{oxonio}\\
        \textbf{2. a)} \textcolor{red}{\ch{Al^3+}}; \textbf{b)} \textcolor{red}{\ch{Pb^4+}}; \textbf{c)} \textcolor{red}{\ch{PH4+}}
    \end{solu}
\end{exercise}

\subsection{Aniones (iones negativos)}

Se forman cuando un átomo o grupo de átomos \negrita{gana} electrones.

\subsubsection{Aniones monoatómicos}

Los forman los \negrita{no metales} (incluido el H) con su número de oxidación negativo. Se nombran con la raíz del nombre del elemento + \blueremark{-uro} (apartado~\ref{sec:raiz}); el oxígeno es la excepción: \negrita{óxido}.
\esquemafig{esq-anion}

El número de carga puede omitirse si no hay ambigüedad.

\ejemplos
\noindent \ballnumber{1} : \ch{Cl-}: cloruro(1-) \(\to\) cloruro\\
\ballnumber{2} : \ch{H-}: hidruro(1-) \(\to\) hidruro\\
\ballnumber{3} : \ch{O^2-}: óxido(2-) \(\to\) óxido\\
\ballnumber{4} : \ch{S^2-}: sulfuro(2-) \(\to\) sulfuro\\
\ballnumber{5} : \ch{N^3-}: nitruro(3-) \(\to\) nitruro

\begin{marker}
    Cuidado con la escritura: \ch{O^2-} es el ion óxido (un átomo con carga \(2-\)), mientras que \ch{O2-} es el ion dióxido(1--) (dos átomos con carga total \(1-\)). Del mismo modo, \ch{N^3-} es el nitruro, pero \ch{N3-} es el ion azida.
\end{marker}

\subsubsection{Aniones con hidrógeno de los grupos 16 y 17}

El \ch{H2S} puede ceder uno o dos hidrones. Si cede los dos queda un anión monoatómico normal; si cede solo uno queda un anión con hidrógeno, que se nombra anteponiendo \blueremark{hidrogeno} (sin tilde):

\ejemplos
\noindent \ballnumber{1} : \ch{H2S -> 2 H+ + S^2-}: ion sulfuro \\
\ballnumber{2} : \ch{H2S -> H+ + HS-}: ion hidrogeno(sulfuro)(1-)\footnote{La forma abreviada «hidrogenosulfuro» no está en la lista de nombres abreviados que acepta la IUPAC. Antiguamente se usaba el prefijo «bi» (\mal{bisulfuro}), nombre totalmente desaconsejado.}

Los oxoaniones (aniones que proceden de los oxoácidos) se estudian después de los oxoácidos, en el apartado~\ref{sec:oxoaniones}.

\begin{exercise}
    \textbf{1.} Nombra los aniones indicando su \concarga: \textbf{a)} \ch{Br-}; \textbf{b)} \ch{S^2-}; \textbf{c)} \ch{N^3-}; \textbf{d)} \ch{O^2-}; \textbf{e)} \ch{HS-}.\\
    \textbf{2.} Escribe la fórmula: \textbf{a)} ion yoduro; \textbf{b)} ion fluoruro; \textbf{c)} ion fosfuro.
    \begin{solu}
        \textbf{1.} \concarga{} \textbf{a)} \textcolor{red}{bromuro(1-)}; \textbf{b)} \textcolor{red}{sulfuro(2-)}; \textbf{c)} \textcolor{red}{nitruro(3-)}; \textbf{d)} \textcolor{red}{óxido(2-)}; \textbf{e)} \textcolor{red}{hidrogeno(sulfuro)(1-)}. \textcolor{blue}{Como no hay ambigüedad, también es correcto omitir la carga: bromuro, sulfuro\dots}\\
        \textbf{2. a)} \textcolor{red}{\ch{I-}}; \textbf{b)} \textcolor{red}{\ch{F-}}; \textbf{c)} \textcolor{blue}{P actúa con \(-3\)} \(\to\) \textcolor{red}{\ch{P^3-}}
    \end{solu}
\end{exercise}

%=====================================================================
\section{Oxoácidos}\label{sec:oxoacidos}
%=====================================================================

Son compuestos ternarios formados por \negrita{hidrógeno}, un elemento central X (normalmente un no metal) y \negrita{oxígeno}. Su fórmula general es \ch{H_aXO_b}. En los oxoácidos el hidrógeno actúa siempre con \(+1\) y el oxígeno con \(-2\).

\subsection{Nombres vulgares admitidos por la IUPAC}

Los oxoácidos se nombran, fundamentalmente, con los \negrita{nombres vulgares admitidos} por la IUPAC (ácido sulfúrico, ácido nítrico\dots). Estos nombres utilizan los prefijos \blueremark{hipo-} y \blueremark{per-} y los sufijos \blueremark{-oso} e \blueremark{-ico} para indicar el número de oxidación del átomo central. Atención: no se trata de una «nomenclatura» que pueda aplicarse a cualquier compuesto; \negrita{solo} son válidos los nombres de la lista.

\esquemafig{esq-oxoacido}

\begin{center}
    \SetTblrInner{rowsep=3pt}
    % --- tabla antigua ---
    % \begin{longtblr}[caption={Oxoácidos con nombre vulgar admitido}]
        % {colspec={Q[c,m] Q[l,m] Q[c,m] Q[l,m]}, row{1}={bg=myOrange3!90!white, fg=white, font=\sffamily}, rowhead=1}
        % \hline
        % Fórmula & Nombre & Nº oxid. de X & Para empezar \\
        % \hline
        % \ch{H2CO3} & ácido carbónico & \(+4\) & \(\star\) \\
        % \ch{HNO3} & ácido nítrico & \(+5\) & \(\star\) \\
        % \ch{H2SO4} & ácido sulfúrico & \(+6\) & \(\star\) \\
        % \ch{H3PO4} & ácido fosfórico & \(+5\) & \(\star\) \\
        % \hline
        % \ch{HNO2} & ácido nitroso & \(+3\) & \\
        % \ch{H2SO3} & ácido sulfuroso & \(+4\) & \\
        % \ch{H3BO3} & ácido bórico & \(+3\) & \\
        % \ch{HClO} & ácido hipocloroso & \(+1\) & \\
        % \ch{HClO2} & ácido cloroso & \(+3\) & \\
        % \ch{HClO3} & ácido clórico & \(+5\) & \\
        % \ch{HClO4} & ácido perclórico & \(+7\) & \\
        % \hline
    % \end{longtblr}
    % --- tabla nueva ---
    \begin{longtable}{@{}clcl@{}}
        \caption{Oxoácidos con nombre vulgar admitido} \\
        \toprule
        \textbf{Fórmula} & \textbf{Nombre} & \textbf{Nº oxid. de X} & \textbf{Para empezar} \\
        \midrule
        \endfirsthead
        \toprule
        \textbf{Fórmula} & \textbf{Nombre} & \textbf{Nº oxid. de X} & \textbf{Para empezar} \\
        \midrule
        \endhead
        \bottomrule
        \endlastfoot
        \ch{H2CO3} & ácido carbónico & \(+4\) & \(\star\) \\[3pt]
        \ch{HNO3} & ácido nítrico & \(+5\) & \(\star\) \\[3pt]
        \ch{H2SO4} & ácido sulfúrico & \(+6\) & \(\star\) \\[3pt]
        \ch{H3PO4} & ácido fosfórico & \(+5\) & \(\star\) \\[3pt]
        \ch{HNO2} & ácido nitroso & \(+3\) &  \\[3pt]
        \ch{H2SO3} & ácido sulfuroso & \(+4\) &  \\[3pt]
        \ch{H3BO3} & ácido bórico & \(+3\) &  \\[3pt]
        \ch{HClO} & ácido hipocloroso & \(+1\) &  \\[3pt]
        \ch{HClO2} & ácido cloroso & \(+3\) &  \\[3pt]
        \ch{HClO3} & ácido clórico & \(+5\) &  \\[3pt]
        \ch{HClO4} & ácido perclórico & \(+7\) &  \\
    \end{longtable}
\end{center}

Igual que los del cloro se nombran los del bromo (\ch{HBrO}, \ch{HBrO2}, \ch{HBrO3}, \ch{HBrO4}) y los del yodo (\ch{HIO}\dots\ \ch{HIO4} ácido peryódico). Los cuatro marcados con \(\star\) son los primeros que conviene aprender.

\begin{marker}
    \negrita{Ácidos del cromo y del manganeso}. La IUPAC \negrita{no} admite nombres vulgares para estos ácidos (\mal{ácido crómico}, \mal{ácido dicrómico}, \mal{ácido permangánico}). Solo acepta los nombres de tres de sus aniones: \ch{CrO4^2-} \negrita{cromato}, \ch{Cr2O7^2-} \negrita{dicromato} y \ch{MnO4-} \negrita{permanganato}. Los ácidos se nombran con la nomenclatura de hidrógeno: \ch{H2CrO4} dihidrogeno(tetraoxidocromato), \ch{H2Cr2O7} dihidrogeno(heptaoxidodicromato), \ch{HMnO4} hidrogeno(tetraoxidomanganato).
\end{marker}

\subsection{Del nombre a la fórmula}

La terminación del nombre nos dice el número de oxidación del átomo central, y este nos da la fórmula. Lo vemos con la «escalera» de los ácidos del cloro:

\figura{fig-escalera}

\begin{center}
    \SetTblrInner{rowsep=4pt}
    % --- tabla antigua ---
    % \begin{tblr}{colspec={Q[l,m] Q[c,m] Q[c,m] Q[c,m] Q[c,m]}, row{1}={bg=myOrange3!90!white, fg=white, font=\sffamily}}
        % \hline
        % Elementos & hipo\dots oso & \dots oso & \dots ico & per\dots ico \\
        % \hline
        % Grupo 17 (Cl, Br, I) & \(+1\) & \(+3\) & \(+5\) & \(+7\)\\
        % Grupo 16 (S) & --- & \(+4\) & \(+6\) & ---\\
        % Grupo 15 (N, P) & --- & \(+3\) & \(+5\) & ---\\
        % Grupo 14 (C) & --- & --- & \(+4\) & ---\\
        % Grupo 13 (B) & --- & --- & \(+3\) & ---\\
        % \hline
    % \end{tblr}
    % --- tabla nueva ---
    \begin{tabular}{@{}lcccc@{}}
        \toprule
        \textbf{Elementos} & \textbf{hipo\dots oso} & \textbf{\dots oso} & \textbf{\dots ico} & \textbf{per\dots ico} \\
        \midrule
        Grupo 17 (Cl, Br, I) & \(+1\) & \(+3\) & \(+5\) & \(+7\) \\[3pt]
        Grupo 16 (S) & --- & \(+4\) & \(+6\) & --- \\[3pt]
        Grupo 15 (N, P) & --- & \(+3\) & \(+5\) & --- \\[3pt]
        Grupo 14 (C) & --- & --- & \(+4\) & --- \\[3pt]
        Grupo 13 (B) & --- & --- & \(+3\) & --- \\
        \bottomrule
    \end{tabular}
\end{center}

\begin{consejos}{oxoácidos}
    \begin{enumerate}[label=\numberballgreen]
        \item \negrita{La terminación -ico} corresponde al número de oxidación igual a la \negrita{última cifra del grupo}: C (grupo 14) \(\to +4\); N (15) \(\to +5\); S (16) \(\to +6\). \negrita{Excepción}: en los halógenos (grupo 17) el \blueremark{-ico} es \(+5\) y el \(+7\) es \blueremark{per-...-ico}.
        \item \negrita{-oso} es dos unidades menos que \blueremark{-ico}, e \negrita{hipo-...-oso} cuatro unidades menos.
        \item \negrita{Número de hidrógenos}: si el número de oxidación de X es \negrita{impar}, el ácido tiene \negrita{1 H}; si es \negrita{par}, tiene \negrita{2 H}. Excepciones: el fósforo y el boro tienen \negrita{3 H} (\ch{H3PO4}, \ch{H3BO3}).
        \item \negrita{Número de oxígenos} (regla del oxígeno):
              \[ \text{nº de O} = \frac{\text{nº de oxidación de X} + \text{nº de H}}{2} \]
              Es simplemente la regla de la suma cero despejada.
    \end{enumerate}
\end{consejos}

\begin{ejemplo}
    \negrita{Ácido sulfuroso}
    \begin{enumerate}[label=\numberballgreen]
        \item El azufre está en el grupo 16, luego \blueremark{-ico} es \(+6\) y \blueremark{-oso} es \(+4\).
        \item \(+4\) es par \(\to\) 2 hidrógenos: \ch{H2SO_x}.
        \item Nº de O \(= \dfrac{4+2}{2} = 3\) \(\to\) \textcolor{red}{\ch{H2SO3}}.
        \item Comprobación (suma cero): \(2\cdot(+1) + (+4) + 3\cdot(-2) = 0\) \checkmark
    \end{enumerate}
\end{ejemplo}

\begin{ejemplo}
    \negrita{Ácido nítrico}
    \begin{enumerate}[label=\numberballgreen]
        \item El nitrógeno está en el grupo 15, luego \blueremark{-ico} es \(+5\).
        \item \(+5\) es impar \(\to\) 1 hidrógeno: \ch{HNO_x}.
        \item Nº de O \(= \dfrac{5+1}{2} = 3\) \(\to\) \textcolor{red}{\ch{HNO3}}.
    \end{enumerate}
\end{ejemplo}

\begin{ejemplo}
    \negrita{Ácido fosfórico}
    \begin{enumerate}[label=\numberballgreen]
        \item El fósforo está en el grupo 15, luego \blueremark{-ico} es \(+5\).
        \item El fósforo es una excepción: lleva 3 hidrógenos: \ch{H3PO_x}.
        \item Nº de O \(= \dfrac{5+3}{2} = 4\) \(\to\) \textcolor{red}{\ch{H3PO4}}.
    \end{enumerate}
\end{ejemplo}

\subsection{De la fórmula al nombre}

\begin{ejemplo}
    \negrita{\ch{HIO4}}
    \begin{enumerate}[label=\numberballgreen]
        \item Calculamos el número de oxidación del yodo (\(x\)) con la suma cero, sabiendo que el H es \(+1\) y el O es \(-2\):
              \begin{equation*}
                  \underbrace{1\cdot(+1)}_{\text{H}} + x + \underbrace{4\cdot(-2)}_{\text{O}} = 0 \quad\Rightarrow\quad x = 8 - 1 = +7
              \end{equation*}
        \item El yodo es un halógeno y \(+7\) es su número de oxidación más alto \(\to\) \blueremark{per-...-ico}.
        \item El nombre es: \textcolor{red}{ácido peryódico}.
    \end{enumerate}
\end{ejemplo}

\begin{exercise}
    \textbf{1.} Formula: \textbf{a)} ácido sulfúrico; \textbf{b)} ácido carbónico; \textbf{c)} ácido nitroso; \textbf{d)} ácido hipobromoso; \textbf{e)} ácido peryódico; \textbf{f)} ácido bórico.\\
    \textbf{2.} Nombra: \textbf{a)} \ch{HClO3}; \textbf{b)} \ch{H2SO3}; \textbf{c)} \ch{HNO3}; \textbf{d)} \ch{HBrO2}; \textbf{e)} \ch{HIO}.
    \begin{solu}
        \textbf{1. a)} \textcolor{blue}{S: -ico \(=+6\), par \(\to\) 2 H; O \(=(6+2)/2=4\)} \(\to\) \textcolor{red}{\ch{H2SO4}}\\
        \textbf{b)} \textcolor{blue}{C: -ico \(=+4\), par \(\to\) 2 H; O \(=(4+2)/2=3\)} \(\to\) \textcolor{red}{\ch{H2CO3}}\\
        \textbf{c)} \textcolor{blue}{N: -oso \(=+3\), impar \(\to\) 1 H; O \(=(3+1)/2=2\)} \(\to\) \textcolor{red}{\ch{HNO2}}\\
        \textbf{d)} \textcolor{blue}{Br: hipo-...-oso \(=+1\); O \(=(1+1)/2=1\)} \(\to\) \textcolor{red}{\ch{HBrO}}\\
        \textbf{e)} \textcolor{blue}{I: per-...-ico \(=+7\); O \(=(7+1)/2=4\)} \(\to\) \textcolor{red}{\ch{HIO4}}\\
        \textbf{f)} \textcolor{blue}{B: \(+3\) y 3 H; O \(=(3+3)/2=3\)} \(\to\) \textcolor{red}{\ch{H3BO3}}\\
        \textbf{2. a)} \textcolor{blue}{Cl \(+5\)} \(\to\) \textcolor{red}{ácido clórico}; \textbf{b)} \textcolor{blue}{S \(+4\)} \(\to\) \textcolor{red}{ácido sulfuroso}; \textbf{c)} \textcolor{blue}{N \(+5\)} \(\to\) \textcolor{red}{ácido nítrico}; \textbf{d)} \textcolor{blue}{Br \(+3\)} \(\to\) \textcolor{red}{ácido bromoso}; \textbf{e)} \textcolor{blue}{I \(+1\)} \(\to\) \textcolor{red}{ácido hipoyodoso}
    \end{solu}
\end{exercise}

\subsection[Otras nomenclaturas de los oxoácidos \ampliacionicono]{Otras nomenclaturas de los oxoácidos \ampliacion}

\begin{enumerate}[label=\numberballgreen]
    \item \textbf{\uline{Nomenclatura de hidrógeno}}. Indica la composición del ácido:
          \esquemafig{esq-hidrogeno}
          \begin{enumerate}[label=\textcolor{red}{\(\bullet\)}]
              \item Prefijo\subin{1}: número de hidrógenos. Prefijo\subin{2}: número de oxígenos. Prefijo\subin{3}: número de átomos centrales (normalmente 1, y se omite).
              \item Todo va junto y lo que va detrás de «hidrogeno» entre paréntesis.
              \item Siempre termina en \blueremark{-ato}: no se usan \blueremark{hipo-}, \blueremark{per-}, \blueremark{-oso} ni \blueremark{-ico}.
              \item \blueremark{hidrogeno} y \blueremark{oxido} van \negrita{sin tilde} porque actúan como prefijos.
          \end{enumerate}
    \item \textbf{\uline{Nomenclatura de adición}}. Escribe la fórmula según su estructura: el átomo central rodeado de grupos \ch{OH} (hidroxido) y átomos de O (oxido), citados por orden alfabético.
\end{enumerate}

\begin{center}
    \SetTblrInner{rowsep=3pt}
    % --- tabla antigua ---
    % \begin{tblr}{colspec={Q[c,m] Q[l,m] Q[c,m] Q[l,m]}, row{1}={bg=myOrange3!90!white, fg=white, font=\sffamily}}
        % \hline
        % Fórmula & Nomenclatura de hidrógeno & Estructura & Nomenclatura de adición \\
        % \hline
        % \ch{HNO3} & hidrogeno(trioxidonitrato) & \ch{[NO2(OH)]} & hidroxidodioxidonitrógeno \\
        % \ch{HClO2} & hidrogeno(dioxidoclorato) & \ch{[ClO(OH)]} & hidroxidooxidocloro \\
        % \ch{HBrO} & hidrogeno(oxidobromato) & \ch{[Br(OH)]} & hidroxidobromo \\
        % \ch{H2SO4} & dihidrogeno(tetraoxidosulfato) & \ch{[SO2(OH)2]} & dihidroxidodioxidoazufre \\
        % \ch{H2CO3} & dihidrogeno(trioxidocarbonato) & \ch{[CO(OH)2]} & dihidroxidooxidocarbono \\
        % \ch{H3PO4} & trihidrogeno(tetraoxidofosfato) & \ch{[PO(OH)3]} & trihidroxidooxidofósforo \\
        % \hline
    % \end{tblr}
    % --- tabla nueva ---
    \begin{tabular}{@{}clcl@{}}
        \toprule
        \textbf{Fórmula} & \textbf{Nomenclatura de hidrógeno} & \textbf{Estructura} & \textbf{Nomenclatura de adición} \\
        \midrule
        \ch{HNO3} & hidrogeno(trioxidonitrato) & \ch{[NO2(OH)]} & hidroxidodioxidonitrógeno \\[3pt]
        \ch{HClO2} & hidrogeno(dioxidoclorato) & \ch{[ClO(OH)]} & hidroxidooxidocloro \\[3pt]
        \ch{HBrO} & hidrogeno(oxidobromato) & \ch{[Br(OH)]} & hidroxidobromo \\[3pt]
        \ch{H2SO4} & dihidrogeno(tetraoxidosulfato) & \ch{[SO2(OH)2]} & dihidroxidodioxidoazufre \\[3pt]
        \ch{H2CO3} & dihidrogeno(trioxidocarbonato) & \ch{[CO(OH)2]} & dihidroxidooxidocarbono \\[3pt]
        \ch{H3PO4} & trihidrogeno(tetraoxidofosfato) & \ch{[PO(OH)3]} & trihidroxidooxidofósforo \\
        \bottomrule
    \end{tabular}
\end{center}

\begin{marker}
    No se debe usar la antigua nomenclatura sistemática: \mal{tetraoxosulfato(VI) de hidrógeno}, \mal{ácido tetraoxosulfúrico(VI)}.
\end{marker}

\begin{exercise}
    \ampliacion\par
    Nombra con la nomenclatura de hidrógeno: \textbf{a)} \ch{HNO2}; \textbf{b)} \ch{H2CO3}; \textbf{c)} \ch{HClO4}; \textbf{d)} \ch{H2SO3}.
    \begin{solu}
        \textbf{a)} \textcolor{red}{hidrogeno(dioxidonitrato)}; \textbf{b)} \textcolor{red}{dihidrogeno(trioxidocarbonato)}; \textbf{c)} \textcolor{red}{hidrogeno(tetraoxidoclorato)}; \textbf{d)} \textcolor{red}{dihidrogeno(trioxidosulfato)}
    \end{solu}
\end{exercise}

%=====================================================================
\section{Oxoaniones}\label{sec:oxoaniones}
%=====================================================================

\subsection{Formación de los oxoaniones}

Todo oxoácido puede ceder iones hidrógeno (hidrones, \ch{H+}). El anión que queda se llama \blueremark{oxoanión} y tiene tantas cargas negativas como hidrones ha perdido.

\begin{center}
    \ch{HNO3 -> H+ + NO3-}\\
    \ch{H2SO4 -> 2 H+ + SO4^2-}\\
    \ch{H3PO4 -> 3 H+ + PO4^3-}
\end{center}

Un ácido con varios hidrógenos puede ceder todos o solo algunos; entonces el anión conserva hidrógeno:
\begin{center}
    \ch{H2CO3 -> H+ + HCO3-}\\
    \ch{H2CO3 -> 2 H+ + CO3^2-}\\
    \ch{H3PO4 -> H+ + H2PO4-}\\
    \ch{H3PO4 -> 2 H+ + HPO4^2-}
\end{center}

En disolución acuosa los hidrones no quedan libres, sino que se unen a moléculas de agua formando iones oxonio:
\begin{center}
    \ch{HNO3 + H2O -> NO3- + H3O+}\\
    \ch{HCl + H2O -> Cl- + H3O+}\\
    \ch{H2CO3 + 2 H2O -> CO3^2- + 2 H3O+}
\end{center}

\subsection{Nombres de los oxoaniones}

\begin{enumerate}[label=\numberballblue]
    \item \textbf{\uline{Nombres vulgares admitidos}}. Se quita la palabra \blueremark{ácido} y se cambian las terminaciones: \blueremark{-oso} \(\to\) \blueremark{-ito} e \blueremark{-ico} \(\to\) \blueremark{-ato}. Los prefijos \blueremark{hipo-} y \blueremark{per-} se mantienen.
          \begin{center}
              \textcolor{myOrange3}{\emph{«El oso toca el pito y Perico el silbato»}}
          \end{center}
          Ojo con el azufre y el fósforo: se dice \blueremark{sulfito}, \blueremark{sulfato}, \blueremark{fosfato}, y \textbf{no} \mal{sulfurito}, \mal{sulfurato} ni \mal{fosforato}.

          \begin{center}
              \SetTblrInner{rowsep=3pt}
              % --- tabla antigua ---
              % \begin{longtblr}[caption={Oxoaniones con nombre vulgar admitido}, note{a}={Estos ácidos no tienen nombre vulgar admitido: se nombran con la nomenclatura de hidrógeno.}]
                  % {colspec={Q[c,m] Q[l,m] Q[c,m] Q[l,m]}, row{1}={bg=myOrange3!90!white, fg=white, font=\sffamily}, rowhead=1}
                  % \hline
                  % Ácido & Nombre del ácido & Anión & Nombre del anión \\
                  % \hline
                  % \ch{HNO3} & ácido nítrico & \ch{NO3-} & nitrato \\
                  % \ch{HNO2} & ácido nitroso & \ch{NO2-} & nitrito \\
                  % \ch{H2SO4} & ácido sulfúrico & \ch{SO4^2-} & sulfato \\
                  % \ch{H2SO3} & ácido sulfuroso & \ch{SO3^2-} & sulfito \\
                  % \ch{H2CO3} & ácido carbónico & \ch{CO3^2-} & carbonato \\
                  % \ch{H3PO4} & ácido fosfórico & \ch{PO4^3-} & fosfato \\
                  % \ch{H3BO3} & ácido bórico & \ch{BO3^3-} & borato \\
                  % \ch{HClO} & ácido hipocloroso & \ch{ClO-} & hipoclorito \\
                  % \ch{HClO2} & ácido cloroso & \ch{ClO2-} & clorito \\
                  % \ch{HClO3} & ácido clórico & \ch{ClO3-} & clorato \\
                  % \ch{HClO4} & ácido perclórico & \ch{ClO4-} & perclorato \\
                  % \ch{H2CrO4} & dihidrogeno(tetraoxidocromato)\TblrNote{a} & \ch{CrO4^2-} & cromato \\
                  % \ch{H2Cr2O7} & dihidrogeno(heptaoxidodicromato)\TblrNote{a} & \ch{Cr2O7^2-} & dicromato \\
                  % \ch{HMnO4} & hidrogeno(tetraoxidomanganato)\TblrNote{a} & \ch{MnO4-} & permanganato \\
                  % \hline
              % \end{longtblr}
              % --- tabla nueva ---
              \begin{longtable}{@{}clcl@{}}
                  \caption{Oxoaniones con nombre vulgar admitido} \\
                  \toprule
                  \textbf{Ácido} & \textbf{Nombre del ácido} & \textbf{Anión} & \textbf{Nombre del anión} \\
                  \midrule
                  \endfirsthead
                  \toprule
                  \textbf{Ácido} & \textbf{Nombre del ácido} & \textbf{Anión} & \textbf{Nombre del anión} \\
                  \midrule
                  \endhead
                  \bottomrule
                  \endlastfoot
                  \ch{HNO3} & ácido nítrico & \ch{NO3-} & nitrato \\[3pt]
                  \ch{HNO2} & ácido nitroso & \ch{NO2-} & nitrito \\[3pt]
                  \ch{H2SO4} & ácido sulfúrico & \ch{SO4^2-} & sulfato \\[3pt]
                  \ch{H2SO3} & ácido sulfuroso & \ch{SO3^2-} & sulfito \\[3pt]
                  \ch{H2CO3} & ácido carbónico & \ch{CO3^2-} & carbonato \\[3pt]
                  \ch{H3PO4} & ácido fosfórico & \ch{PO4^3-} & fosfato \\[3pt]
                  \ch{H3BO3} & ácido bórico & \ch{BO3^3-} & borato \\[3pt]
                  \ch{HClO} & ácido hipocloroso & \ch{ClO-} & hipoclorito \\[3pt]
                  \ch{HClO2} & ácido cloroso & \ch{ClO2-} & clorito \\[3pt]
                  \ch{HClO3} & ácido clórico & \ch{ClO3-} & clorato \\[3pt]
                  \ch{HClO4} & ácido perclórico & \ch{ClO4-} & perclorato \\[3pt]
                  \ch{H2CrO4} & dihidrogeno(tetraoxidocromato)\textsuperscript{a} & \ch{CrO4^2-} & cromato \\[3pt]
                  \ch{H2Cr2O7} & dihidrogeno(heptaoxidodicromato)\textsuperscript{a} & \ch{Cr2O7^2-} & dicromato \\[3pt]
                  \ch{HMnO4} & hidrogeno(tetraoxidomanganato)\textsuperscript{a} & \ch{MnO4-} & permanganato \\
              \end{longtable}
              {\footnotesize
                  \textsuperscript{a}\,Estos ácidos no tienen nombre vulgar admitido: se nombran con la nomenclatura de hidrógeno.\par
              }
          \end{center}

          Si el oxoanión conserva hidrógeno se antepone \blueremark{hidrogeno} (sin tilde), con un prefijo si hay más de uno:

          \ejemplos
          \ballnumber{1} : \ch{HCO3-}: hidrogenocarbonato\footnote{Antiguamente se usaba el prefijo «bi»: \mal{bicarbonato}. Es un nombre muy habitual, pero desaconsejado por la IUPAC.}\\
          \ballnumber{2} : \ch{HSO4-}: hidrogenosulfato\\
          \ballnumber{3} : \ch{HPO4^2-}: hidrogenofosfato\\
          \ballnumber{4} : \ch{H2PO4-}: dihidrogenofosfato

    \item \textbf{\uline{Nomenclatura de adición}} \ampliacion. Como la de hidrógeno de los ácidos, pero sin los hidrógenos ionizados y terminando con el número de carga:
          \esquemafig{esq-oxoanion-adicion}

          \ejemplos
          \ballnumber{1} : \ch{NO3-}: trioxidonitrato(1-) \quad \ballnumber{2} : \ch{SO4^2-}: tetraoxidosulfato(2-)\\
          \ballnumber{3} : \ch{PO4^3-}: tetraoxidofosfato(3-) \quad \ballnumber{4} : \ch{MnO4-}: tetraoxidomanganato(1-)\\
          \ballnumber{5} : \ch{Cr2O7^2-}: heptaoxidodicromato(2-) \quad \ballnumber{6} : \ch{HCO3-}: hidrogeno(trioxidocarbonato)(1-)
\end{enumerate}

\begin{exercise}
    \textbf{1.} Escribe el oxoanión que se obtiene cuando el ácido pierde \emph{todos} sus hidrógenos y nómbralo: \textbf{a)} \ch{HNO3}; \textbf{b)} \ch{H2SO3}; \textbf{c)} \ch{HClO}; \textbf{d)} \ch{HIO4}; \textbf{e)} \ch{H3PO4}.\\
    \textbf{2.} Nombra: \textbf{a)} \ch{ClO2-}; \textbf{b)} \ch{BrO3-}; \textbf{c)} \ch{MnO4-}; \textbf{d)} \ch{CrO4^2-}; \textbf{e)} \ch{H2PO4-}; \textbf{f)} \ch{HCO3-}.
    \begin{solu}
        \textbf{1.} \textcolor{blue}{La carga negativa es igual al número de H perdidos; -oso \(\to\) -ito, -ico \(\to\) -ato.}\\
        \textbf{a)} \textcolor{red}{\ch{NO3-} nitrato}; \textbf{b)} \textcolor{red}{\ch{SO3^2-} sulfito}; \textbf{c)} \textcolor{red}{\ch{ClO-} hipoclorito}; \textbf{d)} \textcolor{red}{\ch{IO4-} peryodato}; \textbf{e)} \textcolor{red}{\ch{PO4^3-} fosfato}\\
        \textbf{2. a)} \textcolor{red}{clorito}; \textbf{b)} \textcolor{red}{bromato}; \textbf{c)} \textcolor{red}{permanganato}; \textbf{d)} \textcolor{red}{cromato}; \textbf{e)} \textcolor{red}{dihidrogenofosfato}; \textbf{f)} \textcolor{red}{hidrogenocarbonato}
    \end{solu}
\end{exercise}

%=====================================================================
\section{Oxosales}\label{sec:oxosales}
%=====================================================================

Son compuestos formados por un \negrita{catión} (metálico o el amonio, \ch{NH4+}) y un \negrita{oxoanión}. Se obtienen al sustituir los hidrógenos de un oxoácido por cationes.

\figura{fig-acido-sal}

\formulan
\begin{enumerate}[label=\numberballgreen]
    \item Escribimos el oxoanión a partir del nombre del ácido del que procede (\blueremark{-ato} \(\to\) \blueremark{-ico}; \blueremark{-ito} \(\to\) \blueremark{-oso}). Su carga es el número de hidrógenos que ha perdido el ácido.
    \item Escribimos el catión con su carga.
    \item Ponemos los subíndices necesarios para que la suma de cargas sea \negrita{cero}. Si hace falta más de un anión poliatómico, va entre paréntesis.
\end{enumerate}

\begin{ejemplo}
    \negrita{Sulfato de aluminio}
    \begin{enumerate}[label=\numberballgreen]
        \item Sulfato \(\to\) ácido sulfúrico \ch{H2SO4} \(\to\) pierde 2 \ch{H+} \(\to\) \ch{SO4^2-}.
        \item Aluminio \(\to\) \ch{Al^3+}.
        \item Suma cero: \(2\cdot(+3) + 3\cdot(-2) = 0\) \(\to\) \textcolor{red}{\ch{Al2(SO4)3}}.
    \end{enumerate}
\end{ejemplo}

\nombran
\begin{enumerate}[label=\numberballblue]
    \item \textbf{\uline{Nombres vulgares admitidos}}:
          \esquemafig{esq-oxosal}

          Para nombrar una sal a partir de la fórmula, la separamos en sus iones. La carga del anión se deduce del número de hidrógenos del ácido o de la carga del catión.

          \ejemplos
          \ballnumber{1} : \ch{Na2SO4 -> 2 Na+ + SO4^2-}: sulfato de sodio \\
          \ballnumber{2} : \ch{CaCO3 -> Ca^2+ + CO3^2-}: carbonato de calcio \\
          \ballnumber{3} : \ch{Ca(IO4)2 -> Ca^2+ + 2 IO4-}: peryodato de calcio \\
          \ballnumber{4} : \ch{CuNO2 -> Cu+ + NO2-}: \conox{} nitrito de cobre(I) \quad \concarga{} nitrito de cobre(1+)\\
          \ballnumber{5} : \ch{Pb(NO3)4 -> Pb^4+ + 4 NO3-}: \conox{} nitrato de plomo(IV) \quad \concarga{} nitrato de plomo(4+)\\
          \ballnumber{6} : \ch{NaClO4 -> Na+ + ClO4-}: perclorato de sodio\\
          \ballnumber{7} : \ch{KMnO4 -> K+ + MnO4-}: permanganato de potasio\\
          \ballnumber{8} : \ch{K2Cr2O7 -> 2 K+ + Cr2O7^2-}: dicromato de potasio

    \item \textbf{\uline{Nomenclaturas de composición y de adición}} \ampliacion:
          \begin{center}
              \SetTblrInner{rowsep=3pt}
              % --- tabla antigua ---
              % \begin{tblr}{colspec={Q[c,m] Q[l,m] Q[l,m]}, row{1}={bg=myOrange3!90!white, fg=white, font=\sffamily}}
                  % \hline
                  % Fórmula & Composición & Adición \\
                  % \hline
                  % \ch{NaNO3} & trioxidonitrato de sodio & trioxidonitrato(1-) de sodio \\
                  % \ch{K2CO3} & trioxidocarbonato de dipotasio & trioxidocarbonato(2-) de potasio \\
                  % \ch{Fe2(SO4)3} & tris(tetraoxidosulfato) de dihierro & tetraoxidosulfato(2-) de hierro(3+) \\
                  % \ch{K2Cr2O7} & heptaoxidodicromato de dipotasio & heptaoxidodicromato(2-) de potasio \\
                  % \ch{Ca3(PO4)2} & bis(tetraoxidofosfato) de tricalcio & tetraoxidofosfato(3-) de calcio \\
                  % \hline
              % \end{tblr}
              % --- tabla nueva ---
              \begin{tabular}{@{}cll@{}}
                  \toprule
                  \textbf{Fórmula} & \textbf{Composición} & \textbf{Adición} \\
                  \midrule
                  \ch{NaNO3} & trioxidonitrato de sodio & trioxidonitrato(1-) de sodio \\[3pt]
                  \ch{K2CO3} & trioxidocarbonato de dipotasio & trioxidocarbonato(2-) de potasio \\[3pt]
                  \ch{Fe2(SO4)3} & tris(tetraoxidosulfato) de dihierro & tetraoxidosulfato(2-) de hierro(3+) \\[3pt]
                  \ch{K2Cr2O7} & heptaoxidodicromato de dipotasio & heptaoxidodicromato(2-) de potasio \\[3pt]
                  \ch{Ca3(PO4)2} & bis(tetraoxidofosfato) de tricalcio & tetraoxidofosfato(3-) de calcio \\
                  \bottomrule
              \end{tabular}
          \end{center}
\end{enumerate}

\begin{exercise}
    \textbf{1.} Formula: \textbf{a)} nitrato de potasio; \textbf{b)} carbonato de calcio; \textbf{c)} hipoclorito de sodio; \textbf{d)} fosfato de magnesio; \textbf{e)} nitrato de hierro(III); \textbf{f)} permanganato de potasio.\\
    \textbf{2.} Nombra indicando, cuando el metal tenga varios, su \conox: \textbf{a)} \ch{Na2SO3}; \textbf{b)} \ch{LiClO4}; \textbf{c)} \ch{CuSO4}; \textbf{d)} \ch{Ca(NO2)2}; \textbf{e)} \ch{Ag2CO3}; \textbf{f)} \ch{Pb(SO4)2}.\\
    \textbf{3.} \ampliacion{} Nombra con \concarga: \textbf{a)} \ch{CuSO4}; \textbf{b)} \ch{Pb(SO4)2}.
    \begin{solu}
        \textbf{1. a)} \textcolor{blue}{\ch{K+} y \ch{NO3-}} \(\to\) \textcolor{red}{\ch{KNO3}}\\
        \textbf{b)} \textcolor{blue}{\ch{Ca^2+} y \ch{CO3^2-}} \(\to\) \textcolor{red}{\ch{CaCO3}}\\
        \textbf{c)} \textcolor{blue}{\ch{Na+} y \ch{ClO-}} \(\to\) \textcolor{red}{\ch{NaClO}}\\
        \textbf{d)} \textcolor{blue}{\ch{Mg^2+} y \ch{PO4^3-}: \(3\cdot(+2)+2\cdot(-3)=0\)} \(\to\) \textcolor{red}{\ch{Mg3(PO4)2}}\\
        \textbf{e)} \textcolor{blue}{\ch{Fe^3+} y \ch{NO3-}} \(\to\) \textcolor{red}{\ch{Fe(NO3)3}}\\
        \textbf{f)} \textcolor{red}{\ch{KMnO4}}\\
        \textbf{2.} \conox{} \textbf{a)} \textcolor{red}{sulfito de sodio}; \textbf{b)} \textcolor{red}{perclorato de litio}; \textbf{c)} \textcolor{red}{sulfato de cobre(II)}; \textbf{d)} \textcolor{red}{nitrito de calcio}; \textbf{e)} \textcolor{red}{carbonato de plata}; \textbf{f)} \textcolor{blue}{2 \ch{SO4^2-} \(\to\) Pb \(+4\)}: \textcolor{red}{sulfato de plomo(IV)}\\
        \textbf{3.} \concarga{} \textbf{a)} \textcolor{red}{sulfato de cobre(2+)}; \textbf{b)} \textcolor{red}{sulfato de plomo(4+)}
    \end{solu}
\end{exercise}

\subsection[Sales ácidas \ampliacionicono]{Sales ácidas \ampliacion}

Proceden de ácidos que han perdido solo \negrita{parte} de sus hidrógenos. Se nombran con el nombre del anión con hidrógeno:

\ejemplos
\ballnumber{1} : \ch{NaHCO3}: hidrogenocarbonato de sodio \\
\ballnumber{2} : \ch{KHSO4}: hidrogenosulfato de potasio \\
\ballnumber{3} : \ch{NH4H2PO4}: dihidrogenofosfato de amonio \\
\ballnumber{4} : \ch{Fe(HCO3)2}: \conox{} hidrogenocarbonato de hierro(II) \quad \concarga{} hidrogenocarbonato de hierro(2+)

\begin{exercise}
    \ampliacion\par
    \textbf{1.} Nombra: \textbf{a)} \ch{KHCO3}; \textbf{b)} \ch{NaHSO4}; \textbf{c)} \ch{Ca(HCO3)2}; \textbf{d)} \ch{NaH2PO4}.\\
    \textbf{2.} Formula: \textbf{a)} hidrogenosulfato de litio; \textbf{b)} hidrogenofosfato de potasio.
    \begin{solu}
        \textbf{1. a)} \textcolor{red}{hidrogenocarbonato de potasio}; \textbf{b)} \textcolor{red}{hidrogenosulfato de sodio}; \textbf{c)} \textcolor{red}{hidrogenocarbonato de calcio}; \textbf{d)} \textcolor{red}{dihidrogenofosfato de sodio}\\
        \textbf{2. a)} \textcolor{blue}{\ch{Li+} y \ch{HSO4-}} \(\to\) \textcolor{red}{\ch{LiHSO4}}; \textbf{b)} \textcolor{blue}{\ch{K+} y \ch{HPO4^2-}} \(\to\) \textcolor{red}{\ch{K2HPO4}}
    \end{solu}
\end{exercise}

\subsection[Sales hidratadas \ampliacionicono]{Sales hidratadas \ampliacion}

Contienen moléculas de agua en su estructura cristalina, que se escriben detrás de un punto. Se nombran con el nombre de la sal, un guion largo, la palabra \negrita{agua} y la proporción entre paréntesis (sal/agua); también es habitual usar el prefijo con \blueremark{-hidrato}:

\ejemplos
\ballnumber{1} : \ch{CuSO4 * 5 H2O}: sulfato de cobre(II)---agua (1/5) o sulfato de cobre(II) pentahidrato\\
\ballnumber{2} : \ch{CaSO4 * 2 H2O}: sulfato de calcio---agua (1/2) o sulfato de calcio dihidrato

\begin{exercise}
    \ampliacion\par
    \textbf{1.} Nombra: \ch{CoCl2 * 6 H2O}.\\
    \textbf{2.} Formula: \textbf{a)} sulfato de magnesio heptahidrato; \textbf{b)} carbonato de sodio---agua (1/10).
    \begin{solu}
        \textbf{1.} \textcolor{red}{cloruro de cobalto(II)---agua (1/6)} o cloruro de cobalto(II) hexahidrato\\
        \textbf{2. a)} \textcolor{red}{\ch{MgSO4 * 7 H2O}}; \textbf{b)} \textcolor{red}{\ch{Na2CO3 * 10 H2O}}
    \end{solu}
\end{exercise}

%=====================================================================
\section{Errores frecuentes}\label{sec:errores}
%=====================================================================

\begin{center}
    \SetTblrInner{rowsep=3pt}
    % --- tabla antigua ---
    % \begin{longtblr}[caption={Nombres y fórmulas incorrectos o desaconsejados}]
        % {colspec={Q[c,m] X[l,m] X[l,m]}, row{1}={bg=red!75!black, fg=white, font=\sffamily}, rowhead=1}
        % \hline
        % Sustancia & Incorrecto / desaconsejado & Correcto \\
        % \hline
        % \ch{SO3} & anhídrido sulfúrico & trióxido de azufre; óxido de azufre(VI) \\
        % \ch{FeO} & óxido ferroso & monóxido de hierro; óxido de hierro(II) \\
        % \ch{Na2O} & monóxido de disodio; óxido de sodio(I) & óxido de disodio; óxido de sodio \\
        % \ch{OCl2} & óxido de dicloro (\ch{Cl2O}) & dicloruro de oxígeno (\ch{OCl2}) \\
        % \ch{N2O4} & simplificar a \ch{NO2} & tetraóxido de dinitrógeno (\ch{N2O4}) \\
        % \ch{Fe2O3} & óxido de hierro (III) (con espacio) & óxido de hierro(III) \\
        % \ch{PH3} & fosfina & fosfano; trihidruro de fósforo \\
        % \ch{HCl} & ácido clorhídrico & cloruro de hidrógeno (ácido clorhídrico es el nombre tradicional de su disolución, no admitido por la IUPAC) \\
        % \ch{H2CrO4} & ácido crómico & dihidrogeno(tetraoxidocromato) \\
        % \ch{HS-} & hidrogenosulfuro & hidrogeno(sulfuro)(1-) \\
        % \ch{H3O+} & hidronio & oxonio \\
        % \ch{Ca^2+} & \ch{Ca^{+2}}; calcio(+2) & \ch{Ca^2+}; calcio(2+) \\
        % \ch{HCO3-} & bicarbonato & hidrogenocarbonato \\
        % \ch{H2SO4} & tetraoxosulfato(VI) de hidrógeno & ácido sulfúrico; dihidrogeno(tetraoxidosulfato) \\
        % \ch{SO4^2-} & sulfurato & sulfato \\
        % \hline
    % \end{longtblr}
    % --- tabla nueva ---
    \begin{longtable}{@{}c L{0.38\linewidth} L{0.38\linewidth}@{}}
        \caption{Nombres y fórmulas incorrectos o desaconsejados} \\
        \toprule
        \textbf{Sustancia} & \textbf{Incorrecto / desaconsejado} & \textbf{Correcto} \\
        \midrule
        \endfirsthead
        \toprule
        \textbf{Sustancia} & \textbf{Incorrecto / desaconsejado} & \textbf{Correcto} \\
        \midrule
        \endhead
        \bottomrule
        \endlastfoot
        \ch{SO3} & anhídrido sulfúrico & trióxido de azufre; óxido de azufre(VI) \\[3pt]
        \ch{FeO} & óxido ferroso & monóxido de hierro; óxido de hierro(II) \\[3pt]
        \ch{Na2O} & monóxido de disodio; óxido de sodio(I) & óxido de disodio; óxido de sodio \\[3pt]
        \ch{OCl2} & óxido de dicloro (\ch{Cl2O}) & dicloruro de oxígeno (\ch{OCl2}) \\[3pt]
        \ch{N2O4} & simplificar a \ch{NO2} & tetraóxido de dinitrógeno (\ch{N2O4}) \\[3pt]
        \ch{Fe2O3} & óxido de hierro (III) (con espacio) & óxido de hierro(III) \\[3pt]
        \ch{PH3} & fosfina & fosfano; trihidruro de fósforo \\[3pt]
        \ch{HCl} & ácido clorhídrico & cloruro de hidrógeno (ácido clorhídrico es el nombre tradicional de su disolución, no admitido por la IUPAC) \\[3pt]
        \ch{H2CrO4} & ácido crómico & dihidrogeno(tetraoxidocromato) \\[3pt]
        \ch{HS-} & hidrogenosulfuro & hidrogeno(sulfuro)(1-) \\[3pt]
        \ch{H3O+} & hidronio & oxonio \\[3pt]
        \ch{Ca^2+} & \ch{Ca^{+2}}; calcio(+2) & \ch{Ca^2+}; calcio(2+) \\[3pt]
        \ch{HCO3-} & bicarbonato & hidrogenocarbonato \\[3pt]
        \ch{H2SO4} & tetraoxosulfato(VI) de hidrógeno & ácido sulfúrico; dihidrogeno(tetraoxidosulfato) \\[3pt]
        \ch{SO4^2-} & sulfurato & sulfato \\
    \end{longtable}
\end{center}

%=====================================================================
\section{Resumen: nombres vulgares admitidos por la IUPAC}\label{sec:vulgares}
%=====================================================================

Se ordenan según el momento en que conviene introducirlos (programación en espiral).

\begin{center}
    \SetTblrInner{rowsep=3pt}
    % --- tabla antigua ---
    % \begin{longtblr}[caption={Nombres vulgares admitidos}]
        % {colspec={Q[c,m] Q[l,m] Q[c,m] Q[l,m]}, row{1}={bg=myOrange3!90!white, fg=white, font=\sffamily}, rowhead=1}
        % \hline
        % Fórmula & Nombre & Fórmula & Nombre \\
        % \hline
        % \SetRow{myOrange3!20!white}\SetCell[c=4]{c} Compuestos binarios\\
        % \hline
        % \ch{H2O} & agua & \ch{NH3} & amoniaco \\
        % \ch{CH4} & metano & \ch{H2O2} & agua oxigenada (peróxido de hidrógeno) \\
        % \ch{SiH4} & silano & \ch{PH3} & fosfano \\
        % \ch{AsH3} & arsano & \ch{SbH3} & estibano \\
        % \hline
        % \SetRow{myOrange3!20!white}\SetCell[c=4]{c} Cationes\\
        % \hline
        % \ch{NH4+} & ion amonio & \ch{H3O+} & ion oxonio \\
        % \hline
        % \SetRow{myOrange3!20!white}\SetCell[c=4]{c} Ácidos y sus aniones\\
        % \hline
        % \ch{HNO3} & ácido nítrico & \ch{NO3-} & nitrato \\
        % \ch{H2SO4} & ácido sulfúrico & \ch{SO4^2-} & sulfato \\
        % \ch{H2CO3} & ácido carbónico & \ch{CO3^2-} & carbonato \\
        % \ch{H3PO4} & ácido fosfórico & \ch{PO4^3-} & fosfato \\
        % \ch{HNO2} & ácido nitroso & \ch{NO2-} & nitrito \\
        % \ch{H2SO3} & ácido sulfuroso & \ch{SO3^2-} & sulfito \\
        % \ch{H3BO3} & ácido bórico & \ch{BO3^3-} & borato \\
        % \ch{HClO} & ácido hipocloroso & \ch{ClO-} & hipoclorito \\
        % \ch{HClO2} & ácido cloroso & \ch{ClO2-} & clorito \\
        % \ch{HClO3} & ácido clórico & \ch{ClO3-} & clorato \\
        % \ch{HClO4} & ácido perclórico & \ch{ClO4-} & perclorato \\
        % \SetRow{myOrange3!20!white}\SetCell[c=4]{c} Aniones con nombre vulgar cuyos ácidos no lo tienen\\
        % \hline
        % \ch{CrO4^2-} & cromato & \ch{Cr2O7^2-} & dicromato \\
        % \ch{MnO4-} & permanganato & & \\
        % \hline
    % \end{longtblr}
    % --- tabla nueva ---
    \begin{longtable}{@{}clcl@{}}
        \caption{Nombres vulgares admitidos} \\
        \toprule
        \textbf{Fórmula} & \textbf{Nombre} & \textbf{Fórmula} & \textbf{Nombre} \\
        \midrule
        \endfirsthead
        \toprule
        \textbf{Fórmula} & \textbf{Nombre} & \textbf{Fórmula} & \textbf{Nombre} \\
        \midrule
        \endhead
        \bottomrule
        \endlastfoot
        \multicolumn{4}{c}{Compuestos binarios} \\[3pt]
        \ch{H2O} & agua & \ch{NH3} & amoniaco \\[3pt]
        \ch{CH4} & metano & \ch{H2O2} & agua oxigenada (peróxido de hidrógeno) \\[3pt]
        \ch{SiH4} & silano & \ch{PH3} & fosfano \\[3pt]
        \ch{AsH3} & arsano & \ch{SbH3} & estibano \\[3pt]
        \multicolumn{4}{c}{Cationes} \\[3pt]
        \ch{NH4+} & ion amonio & \ch{H3O+} & ion oxonio \\[3pt]
        \multicolumn{4}{c}{Ácidos y sus aniones} \\[3pt]
        \ch{HNO3} & ácido nítrico & \ch{NO3-} & nitrato \\[3pt]
        \ch{H2SO4} & ácido sulfúrico & \ch{SO4^2-} & sulfato \\[3pt]
        \ch{H2CO3} & ácido carbónico & \ch{CO3^2-} & carbonato \\[3pt]
        \ch{H3PO4} & ácido fosfórico & \ch{PO4^3-} & fosfato \\[3pt]
        \ch{HNO2} & ácido nitroso & \ch{NO2-} & nitrito \\[3pt]
        \ch{H2SO3} & ácido sulfuroso & \ch{SO3^2-} & sulfito \\[3pt]
        \ch{H3BO3} & ácido bórico & \ch{BO3^3-} & borato \\[3pt]
        \ch{HClO} & ácido hipocloroso & \ch{ClO-} & hipoclorito \\[3pt]
        \ch{HClO2} & ácido cloroso & \ch{ClO2-} & clorito \\[3pt]
        \ch{HClO3} & ácido clórico & \ch{ClO3-} & clorato \\[3pt]
        \ch{HClO4} & ácido perclórico & \ch{ClO4-} & perclorato \\[3pt]
        \multicolumn{4}{c}{Aniones con nombre vulgar cuyos ácidos no lo tienen} \\[3pt]
        \ch{CrO4^2-} & cromato & \ch{Cr2O7^2-} & dicromato \\[3pt]
        \ch{MnO4-} & permanganato &  &  \\
    \end{longtable}
\end{center}

%=====================================================================
\section{Comprueba lo aprendido}\label{sec:practica}
%=====================================================================

\begin{exercise}
    \negrita{Nombra} cada compuesto de dos formas, \conpref{} y \conox: \textbf{a)} \ch{CO2}; \textbf{b)} \ch{K2O}; \textbf{c)} \ch{Fe2O3}; \textbf{d)} \ch{CoH2}; \textbf{e)} \ch{PCl5}; \textbf{f)} \ch{AgCl}; \textbf{g)} \ch{PbS2}; \textbf{h)} \ch{Al(OH)3}; \textbf{i)} \ch{SO2}.\\
    \ampliacion{} Después, nombra con \concarga{} los que sean iónicos y tengan un metal con varias cargas posibles.
    \begin{solu}
        \textbf{a)} \conpref{} \textcolor{red}{dióxido de carbono} \quad \conox{} \textcolor{red}{óxido de carbono(IV)}\\
        \textbf{b)} \conpref{} \textcolor{red}{óxido de dipotasio} \quad \conox{} \textcolor{red}{óxido de potasio}\\
        \textbf{c)} \conpref{} \textcolor{red}{trióxido de dihierro} \quad \conox{} \textcolor{red}{óxido de hierro(III)}\\
        \textbf{d)} \conpref{} \textcolor{red}{dihidruro de cobalto} \quad \conox{} \textcolor{red}{hidruro de cobalto(II)}\\
        \textbf{e)} \conpref{} \textcolor{red}{pentacloruro de fósforo} \quad \conox{} \textcolor{red}{cloruro de fósforo(V)}\\
        \textbf{f)} \conpref{} \textcolor{red}{cloruro de plata} \quad \conox{} \textcolor{red}{cloruro de plata}\\
        \textbf{g)} \conpref{} \textcolor{red}{disulfuro de plomo} \quad \conox{} \textcolor{red}{sulfuro de plomo(IV)}\\
        \textbf{h)} \conpref{} \textcolor{red}{trihidróxido de aluminio} \quad \conox{} \textcolor{red}{hidróxido de aluminio}\\
        \textbf{i)} \conpref{} \textcolor{red}{dióxido de azufre} \quad \conox{} \textcolor{red}{óxido de azufre(IV)}\\
        \concarga{} \textbf{c)} \textcolor{red}{óxido de hierro(3+)}; \textbf{d)} \textcolor{red}{hidruro de cobalto(2+)}; \textbf{g)} \textcolor{red}{sulfuro de plomo(4+)}. \textcolor{blue}{\ch{CO2}, \ch{PCl5} y \ch{SO2} son covalentes: no admiten número de carga.}
    \end{solu}
\end{exercise}

\begin{exercise}
    \negrita{Formula} los siguientes compuestos: \textbf{a)} óxido de aluminio; \textbf{b)} cloruro de cinc; \textbf{c)} trihidruro de níquel; \textbf{d)} yoduro de plomo(II); \textbf{e)} peróxido de potasio; \textbf{f)} ácido nítrico; \textbf{g)} sulfato de cobre(II); \textbf{h)} carbonato de sodio; \textbf{i)} hidrogenocarbonato de sodio; \textbf{j)} ion amonio.
    \begin{solu}
        \textbf{a)} \textcolor{red}{\ch{Al2O3}}; \textbf{b)} \textcolor{red}{\ch{ZnCl2}}; \textbf{c)} \textcolor{red}{\ch{NiH3}}; \textbf{d)} \textcolor{red}{\ch{PbI2}}; \textbf{e)} \textcolor{red}{\ch{K2O2}}; \textbf{f)} \textcolor{red}{\ch{HNO3}}; \textbf{g)} \textcolor{red}{\ch{CuSO4}}; \textbf{h)} \textcolor{red}{\ch{Na2CO3}}; \textbf{i)} \textcolor{red}{\ch{NaHCO3}}; \textbf{j)} \textcolor{red}{\ch{NH4+}}
    \end{solu}
\end{exercise}

% ---------------------------------------------------------------------
% Soluciones de todos los ejercicios
% ---------------------------------------------------------------------
\Closesolutionfile{ans}
\newpage
\section*{Soluciones de los ejercicios}
\addcontentsline{toc}{section}{Soluciones de los ejercicios}
\input{ans1}

\end{document}
